% Chapter 3. All quoted numbers come from chapters/ch03-numbers.tex, written by code/ch03/run_all.py.
% Generated by code/ch03/run_all.py. Do not edit by hand.
\newcommand{\bgkBeta}{0.5826}
\newcommand{\bmMaxBGKBig}{30.86}
\newcommand{\bmMaxBGKSixteen}{25.19}
\newcommand{\bmMaxCont}{31.73}
\newcommand{\bmMaxMCBig}{30.50}
\newcommand{\bmMaxMCSixteen}{25.60}
\newcommand{\bmMaxSeBig}{0.23}
\newcommand{\bmMaxZAt}{1,024}
\newcommand{\bmMaxZmax}{1.54}
\newcommand{\canCarryBp}{366}
\newcommand{\canConvex}{0.084}
\newcommand{\canConvexBp}{8.4}
\newcommand{\canConvexRealMed}{0.069}
\newcommand{\canConvexShare}{2.2}
\newcommand{\canEnd}{September 2026}
\newcommand{\canEndEmp}{27.0}
\newcommand{\canEndModel}{27.1}
\newcommand{\canHitDn}{52.5}
\newcommand{\canHitEitherEmp}{87.6}
\newcommand{\canHitEitherModel}{95.4}
\newcommand{\canHitEmp}{49.4}
\newcommand{\canHitModel}{54.2}
\newcommand{\canHitMonthEmp}{42.4}
\newcommand{\canHitRatio}{1.83}
\newcommand{\canHitUp}{46.3}
\newcommand{\canN}{6,437}
\newcommand{\canQvLinDev}{23}
\newcommand{\canQvMaxShare}{7}
\newcommand{\canQvMaxSq}{0.0400}
\newcommand{\canQvTwentyFive}{0.555}
\newcommand{\canRvDailyMed}{0.75}
\newcommand{\canRvMax}{1.29}
\newcommand{\canRvMaxYear}{2023}
\newcommand{\canRvMin}{0.50}
\newcommand{\canRvMinYear}{2018}
\newcommand{\canRvMonthlyMed}{0.70}
\newcommand{\canRvNoiseHi}{1.24}
\newcommand{\canRvNoiseLo}{0.72}
\newcommand{\canRvRatioHi}{1.14}
\newcommand{\canRvRatioLo}{0.74}
\newcommand{\canRvRatioMed}{0.96}
\newcommand{\canRvWeeklyMed}{0.77}
\newcommand{\canSigma}{0.82}
\newcommand{\canStart}{January 2001}
\newcommand{\canStarts}{6,185}
\newcommand{\canYearsIndep}{25}
\newcommand{\canYieldEnd}{3.69}
\newcommand{\canYieldEndCc}{3.66}
\newcommand{\canYieldEndDate}{29 September 2026}
\newcommand{\caseADF}{\ensuremath{-}1.17}
\newcommand{\caseADFcrit}{\ensuremath{-}2.88}
\newcommand{\caseADFp}{0.69}
\newcommand{\caseAmp}{52}
\newcommand{\caseBand}{2.23}
\newcommand{\caseBiasRatio}{2.23}
\newcommand{\caseBoot}{2,000}
\newcommand{\caseBootValid}{99.9}
\newcommand{\caseEnd}{September 2026}
\newcommand{\caseHalf}{3.03}
\newcommand{\caseHalfHi}{4.66}
\newcommand{\caseHalfLo}{0.63}
\newcommand{\caseKappa}{0.229}
\newcommand{\caseKappaBias}{0.282}
\newcommand{\caseKappaBootMean}{0.510}
\newcommand{\caseKappaKendall}{\ensuremath{-}0.000}
\newcommand{\caseLast}{3.69}
\newcommand{\caseLongADF}{\ensuremath{-}2.06}
\newcommand{\caseLongADFp}{0.26}
\newcommand{\caseLongHalf}{2.78}
\newcommand{\caseLongKappa}{0.249}
\newcommand{\caseLongN}{309}
\newcommand{\caseLongStart}{January 2001}
\newcommand{\caseLongTheta}{2.34}
\newcommand{\caseN}{210}
\newcommand{\caseNK}{209}
\newcommand{\casePhi}{0.9811}
\newcommand{\casePhiKendall}{1.0000}
\newcommand{\casePhiKendallAlt}{0.9999}
\newcommand{\caseSdInf}{1.14}
\newcommand{\caseSigma}{0.770}
\newcommand{\caseStart}{April 2009}
\newcommand{\caseTheta}{2.34}
\newcommand{\exEight}{3}
\newcommand{\exEightB}{1.62}
\newcommand{\exFive}{9}
\newcommand{\exFourAcf}{0.6065}
\newcommand{\exFourHalf}{1.386}
\newcommand{\exFourSd}{2.00}
\newcommand{\exOneCov}{1}
\newcommand{\exOneVar}{5}
\newcommand{\exSeven}{1}
\newcommand{\exSixP}{74.57}
\newcommand{\exSixPZero}{66.67}
\newcommand{\exSixTau}{0.4743}
\newcommand{\exSixTauZero}{0.50}
\newcommand{\exThreeLoss}{39.1}
\newcommand{\exThreeMean}{117.35}
\newcommand{\exThreeMedian}{110.24}
\newcommand{\exTwo}{12}
\newcommand{\exactMeanOne}{108.3287}
\newcommand{\exitPaths}{8,000}
\newcommand{\exitPup}{62.25}
\newcommand{\exitPupMC}{61.70}
\newcommand{\exitTau}{0.2449}
\newcommand{\exitTauMC}{0.2446}
\newcommand{\exitTauZero}{0.2500}
\newcommand{\exitTauZeroMC}{0.2502}
\newcommand{\fkBS}{9.4134}
\newcommand{\fkErrCallBig}{0.00004}
\newcommand{\fkMC}{9.4063}
\newcommand{\fkMCse}{0.0105}
\newcommand{\fkOUExact}{0.14065}
\newcommand{\fkOUMC}{0.1404}
\newcommand{\fkOUMCse}{0.0004}
\newcommand{\fkOUPDE}{0.14065}
\newcommand{\fkPDE}{9.4134}
\newcommand{\fkPaths}{1,000,000}
\newcommand{\fkSlopeCall}{\ensuremath{-}1.95}
\newcommand{\fkSlopeOU}{\ensuremath{-}2.00}
\newcommand{\gbmGrowthHigh}{\ensuremath{-}4.500}
\newcommand{\gbmGrowthLow}{4.875}
\newcommand{\gbmMeanHigh}{1,102}
\newcommand{\gbmMeanLow}{1,102}
\newcommand{\gbmMeanLowTen}{222.6}
\newcommand{\gbmMedianHigh}{25.9}
\newcommand{\gbmMedianLow}{431.7}
\newcommand{\gbmMedianLowTen}{162.8}
\newcommand{\gbmProbLossHigh}{68.9}
\newcommand{\gbmProbLossLow}{14.3}
\newcommand{\girDirect}{0.3519}
\newcommand{\girDirectSE}{0.0009}
\newcommand{\girExact}{0.3509}
\newcommand{\girPaths}{400,000}
\newcommand{\girWeighted}{0.3539}
\newcommand{\girWeightedOdds}{27}
\newcommand{\girWeightedOddsOne}{54}
\newcommand{\girWeightedSE}{0.0014}
\newcommand{\girWeightedZ}{2.1}
\newcommand{\isBig}{13.51}
\newcommand{\isBigSE}{0.045}
\newcommand{\isP}{13.50}
\newcommand{\isPlainBig}{13.97}
\newcommand{\isPlainBigSE}{0.68}
\newcommand{\isPlainSmall}{0.0}
\newcommand{\isPlainSmallSE}{0.0}
\newcommand{\isPlainSmallTrueSE}{11.6}
\newcommand{\isRatio}{218}
\newcommand{\isSmall}{14.01}
\newcommand{\isSmallSE}{0.81}
\newcommand{\isoEExact}{3.195}
\newcommand{\isoESE}{0.315}
\newcommand{\isoEVar}{3.349}
\newcommand{\isoTExact}{0.3333}
\newcommand{\isoTVar}{0.3293}
\newcommand{\itoBigN}{65,536}
\newcommand{\itoLimit}{\ensuremath{-}0.1072}
\newcommand{\itoMean}{\ensuremath{-}0.0016}
\newcommand{\itoPaths}{20,000}
\newcommand{\itoSmallN}{16}
\newcommand{\itoSumBig}{\ensuremath{-}0.1046}
\newcommand{\itoSumSmall}{\ensuremath{-}0.1351}
\newcommand{\itoVar}{0.5043}
\newcommand{\itoWT}{\ensuremath{-}0.8864}
\newcommand{\lemmaN}{4,096}
\newcommand{\lemmaPaths}{4,000}
\newcommand{\lemmaRmsCube}{0.0464}
\newcommand{\lemmaRmsExp}{0.0196}
\newcommand{\lemmaSlopeCube}{\ensuremath{-}0.50}
\newcommand{\lemmaSlopeExp}{\ensuremath{-}0.50}
\newcommand{\levyCovQ}{0.250}
\newcommand{\levyIncCov}{0.003}
\newcommand{\levyPaths}{40,000}
\newcommand{\levyVarHalf}{0.500}
\newcommand{\levyVarOne}{1.008}
\newcommand{\mgExp}{1.0029}
\newcommand{\mgExpSE}{0.0029}
\newcommand{\mgExpThree}{0.985}
\newcommand{\mgExpThreeSD}{90.0}
\newcommand{\mgExpThreeSE}{0.077}
\newcommand{\mgExpThreeShare}{6.7}
\newcommand{\mgExpThreeTrueSE}{0.201}
\newcommand{\mgPaths}{200,000}
\newcommand{\mgRejShare}{18}
\newcommand{\mgReps}{400}
\newcommand{\mgSEHalfShare}{75}
\newcommand{\mgSERatio}{2.6}
\newcommand{\mgSq}{0.0019}
\newcommand{\mgSqSE}{0.0032}
\newcommand{\naiveGapCorr}{0.99967}
\newcommand{\naiveGapRms}{0.0449}
\newcommand{\naiveRmsCube}{1.713}
\newcommand{\naiveRmsExp}{0.796}
\newcommand{\normCdfOne}{0.8413}
\newcommand{\ouAmp}{10}
\newcommand{\ouHalfLife}{0.578}
\newcommand{\ouKappaBiasRel}{7}
\newcommand{\ouKappaSd}{0.24}
\newcommand{\ouKendall}{\ensuremath{-}0.0062}
\newcommand{\ouKendallKappa}{0.08}
\newcommand{\ouMeanOne}{5.11}
\newcommand{\ouPhiBias}{\ensuremath{-}0.0063}
\newcommand{\ouPhiBiasRel}{0.7}
\newcommand{\ouPhiHatMean}{0.8985}
\newcommand{\ouPhiMonthly}{0.9048}
\newcommand{\ouProbNeg}{1.01}
\newcommand{\ouSdInf}{1.29}
\newcommand{\qvMSEratio}{0.021}
\newcommand{\qvPathEight}{0.9501}
\newcommand{\qvPathFour}{0.7334}
\newcommand{\qvPathSixteen}{0.9889}
\newcommand{\qvPathTwelve}{0.9570}
\newcommand{\qvSlope}{\ensuremath{-}0.499}
\newcommand{\reflA}{4}
\newcommand{\reflExact}{38.33}
\newcommand{\reflMC}{38.39}
\newcommand{\reflMCse}{0.08}
\newcommand{\reflN}{20}
\newcommand{\reflNaive}{25.17}
\newcommand{\rwCdfHundred}{0.8644}
\newcommand{\rwCdfTen}{0.8281}
\newcommand{\rwCdfThousand}{0.8365}
\newcommand{\rwCovCheck}{103.3}
\newcommand{\rwKolSlope}{\ensuremath{-}0.498}
\newcommand{\rwVarFourHundred}{403.8}
\newcommand{\snipExcess}{0.115}
\newcommand{\snipExcessRest}{0.03}
\newcommand{\snipHalf}{0.527}
\newcommand{\snipKappa}{1.315}
\newcommand{\stratGapMean}{0.4998}
\newcommand{\stratGapSD}{0.0389}
\newcommand{\stratLimit}{0.3928}
\newcommand{\stratSumBig}{0.3901}
\newcommand{\stratSumSmall}{0.3685}
\newcommand{\strongEuler}{0.49}
\newcommand{\strongEulerSixteen}{2.388}
\newcommand{\strongMilSixteen}{0.206}
\newcommand{\strongMilstein}{0.95}
\newcommand{\strongPaths}{4,000}
\newcommand{\tvMeanSixteen}{204.3}
\newcommand{\tvPathEight}{12.5}
\newcommand{\tvPathFour}{2.4}
\newcommand{\tvPathSixteen}{202.9}
\newcommand{\tvPathTwelve}{49.8}
\newcommand{\weakLead}{\ensuremath{-}0.0289}
\newcommand{\weakMeanTwelve}{108.3000}
\newcommand{\weakSlopeMean}{0.997}
\newcommand{\weakSlopeSecond}{0.987}

\begin{refsection}
\setcounter{chapter}{2}
\chapter{Random Walks and Stochastic Calculus}
\label{ch:stochastic}

Prices, rates and appraised values move a little every day, and nobody can say in advance which way.
The five year Government of Canada yield, which runs through this chapter's examples, has had a
realised volatility of about \canRvDailyMed{} of a percentage point a year in the median year since
2001, built up from daily changes of a few basis points. A model that respects this has to be random, and for questions about hedging, valuation
or the chance of crossing a covenant level between two observation dates it helps to work in
continuous time. This chapter builds the calculus that such models need. It starts from a
coin-tossing walk and takes a limit to reach Brownian motion. The limit has a property ordinary
calculus never meets: each squared increment vanishes as the step shrinks, but their sum over finer
and finer partitions converges to the elapsed time. Nearly every correction term in this chapter traces back to it.

Once that property is in hand, the random walk leads directly to the continuous time tools used in pricing and risk management.
It tells us how rough a yield path is and why, for a Brownian path, sampling more often sharpens the
volatility estimate without changing what is being estimated
(\cref{sec3:rw,sec3:bm,sec3:qv}), and through the reflection principle how often a yield crosses a
level between the month-end dates at which it is usually recorded. It forces a new integral for trading gains and a chain rule with
an extra term, which is where a bond's convexity earns its keep (\cref{sec3:ito,sec3:lemma}). From
there come stopping rules for barriers (\cref{sec3:mart}), simulation schemes (\cref{sec3:sde}), the
change of measure that moves a drift or makes a rare loss cheap to estimate (\cref{sec3:girsanov}),
and the link between expectations and differential equations that Chapter~5 uses to price options
(\cref{sec3:fk}). The case study then fits a mean-reverting model to the five year yield and finds
that the data say less than the fitted numbers suggest (\cref{sec3:case}).

\begin{keybox}[How to read this chapter]
Proofs are given in full where they are short and in outline, with the gap named, where a complete
proof needs measure theory beyond this book; \textcite{karatzas1998} and \textcite{morters2010} fill
those gaps. The figures and calculations, apart from values attributed to published data, are generated by eight
independently seeded Python scripts in \texttt{code/ch03/}, checked against closed forms and
supplemented by independent R and GNU Octave calculations. The notes at the end of the chapter give the
dependencies, seeds and verification details.
\end{keybox}
% =====================================================================================
\section{Why a monthly model is too coarse}
\label{sec3:why}

A discrete model with one step a month is often good enough to describe a fund. Monthly steps may be
too coarse to assess the hedge error and transaction costs of an option that is rebalanced more
often, because the hedge has to be adjusted as the price moves, and the cost of hedging
depends on how much the price moves between adjustments. Shrinking the step does not make the problem
go away. It changes it into a question about limits, and the limit turns out to have rules of its own.

Alternative assets give several reasons to care. Commodity spreads, credit spreads and short rates
are often modelled as mean reverting, although the strength and stability of the reversion must be
tested, and the case study tests it on the Canadian five year yield. If the pull exists, its speed helps frame how long a relative value position may need to be held. Barrier events matter to leveraged funds: a margin loan or a mark-to-market trigger acts the first
time a value crosses its level, even if the value has recovered by the next month end, while a
covenant tested only on reporting dates is a barrier monitored on those dates alone. And the valuation of
options, carried interest and guarantees, met in Chapter~2 and again in Chapters~5 and~6, rests on
expectations that can be computed either by simulation or by solving a differential equation. Under conditions the
two agree, and \cref{sec3:fk} proves that a smooth solution of the equation is the expectation.

\FloatBarrier
% =====================================================================================
\section{Random walks and the size of a typical move}
\label{sec3:rw}

\begin{definition}\label{def3:srw}
Let $\xi_1,\xi_2,\dots$ be independent with $P(\xi_k=1)=P(\xi_k=-1)=\tfrac12$. The \emph{simple
random walk} is $S_0=0$ and $S_n=\xi_1+\dots+\xi_n$. Write $\mathcal F_n$ for the information in
$\xi_1,\dots,\xi_n$.
\end{definition}

\begin{proposition}\label{prop3:srw}
For $m\le n$, $\E[S_n]=0$, $\Var[S_n]=n$, $\Cov(S_m,S_n)=m$, and $\E[S_n\mid\mathcal F_m]=S_m$.
\end{proposition}

\begin{proof}
Each $\xi_k$ has mean $0$ and variance $1$, so the first two claims follow from linearity and
independence. For the covariance write $S_n = S_m + (S_n-S_m)$; the second term is a sum of steps after
$m$, independent of $S_m$ with mean zero, so $\Cov(S_m,S_n)=\Var[S_m]=m$. The same split gives
$\E[S_n\mid\mathcal F_m] = S_m + \E[S_n-S_m] = S_m$.
\end{proof}

The last property makes $S$ a \emph{martingale}: the best forecast of a future value, in mean square and given the path so far, is the present
value. The covariance $\min(m,n)$ will reappear unchanged in the continuous limit.

To pass to that limit, speed the walk up and shrink its steps. With $n$ steps per unit of time and step
size $1/\sqrt n$, the rescaled path $t\mapsto S_{\lfloor nt\rfloor}/\sqrt n$ has variance $\lfloor nt\rfloor/n$ at time $t$, within
$1/n$ of $t$ for every $n$. Panel (a) of \cref{fig3:rwscaling} shows three such paths. At $n=16$ the
steps are visible; at $n=4{,}096$ the path looks like a price chart.

\begin{theorem}[Donsker]\label{thm3:donsker}
The rescaled walks $W^{(n)}_t = S_{\lfloor nt\rfloor}/\sqrt n$ converge in distribution, as random
continuous functions on $[0,1]$ after linear interpolation, to a limit $W$ with independent Gaussian
increments $W_t-W_s\sim\N(0,t-s)$.
\end{theorem}

We do not prove this here; \textcite{billingsley1999} gives the proof. Its one-dimensional shadow is the central limit theorem, and for this lattice distribution the largest gap from the normal distribution function can be
computed exactly at each $n$. Panel (b)
plots the largest gap between the exact distribution function of $S_n/\sqrt n$ and the normal one,
computed from the binomial law. On log-log axes the fitted slope over $10\le n\le10{,}000$ is \rwKolSlope, an empirical value consistent
with an error that falls like $1/\sqrt n$. Lattice effects make the rate uneven at a fixed point: $P(S_n/\sqrt n\le 1)$ is
\rwCdfTen{} at $n=10$, \rwCdfHundred{} at $n=100$ and \rwCdfThousand{} at $n=1{,}000$, against
\normCdfOne{} in the limit. At $n=100$ the point $1$ is itself a value the walk can take, and the whole
mass on it is counted.

\begin{figure}[t]
  \centering
  \includegraphics[width=\textwidth]{ch03/fig01_rw_scaling}
  \caption{(a) Random walks with $n$ steps on $[0,1]$, scaled by $1/\sqrt n$. (b) The exact largest
  distance between the distribution function of $S_n/\sqrt n$ and the standard normal one, which falls
  like $1/\sqrt n$.}
  \label{fig3:rwscaling}
\end{figure}

\subsection{The reflection principle}

How likely is the walk to reach a level $a$ at some time up to $n$? The answer needs the whole path,
yet a symmetry reduces it to the endpoint.

\begin{proposition}\label{prop3:srwrefl}
For integers $a\ge 1$ and $n\ge 1$,
\begin{equation}\label{eq3:srwrefl}
  P\Bigl(\max_{k\le n} S_k \ge a\Bigr) = P(S_n\ge a) + P(S_n > a) .
\end{equation}
\end{proposition}

\begin{proof}
Every path of length $n$ has probability $2^{-n}$, so it is enough to count paths. Split the paths whose
maximum reaches $a$ by their endpoint. Those with $S_n\ge a$ reach $a$ trivially and are counted by
$P(S_n\ge a)$. For a path with $S_n<a$ that does reach $a$, let $\tau$ be the first time it does, and
reflect the part after $\tau$ in the level $a$: replace $S_k$ by $2a-S_k$ for $k>\tau$. The new path
ends at $2a-S_n>a$. The map is a bijection between paths that reach $a$ and end below it and paths
that end above $a$, since every path ending above $a$ has a first visit to $a$ and reflecting again
recovers the original. Hence the second group has probability $P(S_n>a)$.
\end{proof}

With $n=\reflN$ and $a=\reflA$ the formula gives \reflExact\%, while the endpoint alone gives only
$P(S_n\ge a)=\reflNaive\%$. Simulating 400{,}000 walks gives \reflMC\% with a standard error of
\reflMCse{} points, and the code also confirms the value by listing all $2^{20}$ paths. Panel (a) of
\cref{fig3:reflection} shows the same reflection applied to a Brownian path.

\FloatBarrier
% =====================================================================================
\section{Brownian motion, built from midpoints}
\label{sec3:bm}

\begin{definition}\label{def3:bm}
A \emph{standard Brownian motion} on a probability space $(\Omega,\mathcal F,P)$ is a process
$W=(W_t)_{t\ge 0}$ with
\begin{enumerate}[label=(\roman*)]
  \item $W_0=0$;
  \item independent increments: for $0\le t_0<t_1<\dots<t_k$, the variables $W_{t_j}-W_{t_{j-1}}$ are independent;
  \item Gaussian increments: $W_t-W_s\sim\N(0,t-s)$ for $s<t$;
  \item continuous paths: with probability one, $t\mapsto W_t(\omega)$ is continuous.
\end{enumerate}
\end{definition}

Properties (i) to (iii) fix all the finite-dimensional laws, and (iv) is what makes the object a
random function rather than a family of random numbers. That all four can hold together is not
obvious. The midpoint construction below proves existence and also provides a simulation algorithm.

\subsection{L\'evy's construction}

The idea is to draw the endpoint first and then fill in midpoints, each time using the right
conditional law. The next lemma supplies that law.

\begin{lemma}\label{lem3:bridge}
Let $s<u$, $m=(s+u)/2$ and $h=u-s$. If $W$ has the finite-dimensional laws of
\cref{def3:bm}, then conditional on $W_s=x$ and $W_u=y$,
\begin{equation}\label{eq3:bridge}
  W_m \sim \N\Bigl(\frac{x+y}{2},\ \frac h4\Bigr),
\end{equation}
and, conditional on both endpoint values, $W_m$ is independent of the path outside $[s,u]$. This is
the conditional independence property of the Brownian bridge.
\end{lemma}

\begin{proof}
By (ii) and (iii), $A=W_m-W_s$ and $B=W_u-W_m$ are independent $\N(0,h/2)$ and independent of the
increments outside $[s,u]$. Their sum $A+B=W_u-W_s$ and difference $A-B$ are jointly Gaussian with
$\Cov(A+B,A-B)=\Var A-\Var B=0$, so they are independent. Hence
$W_m = W_s + \tfrac12(A+B) + \tfrac12(A-B)$, and given $W_s=x$ and $W_u=y$ the first two terms equal
$(x+y)/2$ while $\tfrac12(A-B)\sim\N(0,h/4)$ is untouched by the conditioning. The independence claim
holds because $A-B$ is uncorrelated with every coordinate of the vector made of $W_s$, $A+B$ and the
increments outside $[s,u]$, and the whole collection is jointly Gaussian, so $A-B$ is independent of
that vector as a whole; conditioning on the endpoints therefore leaves it unchanged and independent
of the outside path.
\end{proof}

\begin{theorem}[L\'evy]\label{thm3:levy}
Let $Z$ and $Z_{n,k}$ be independent standard normals. Set $W^{(0)}_0=0$, $W^{(0)}_1=Z$ and, at level
$n\ge1$, keep the values on the grid $\{k2^{-(n-1)}\}$ and define each new midpoint by
\begin{equation}\label{eq3:levy}
  W^{(n)}_{(2k+1)2^{-n}} = \tfrac12\bigl(W^{(n-1)}_{k2^{-(n-1)}} + W^{(n-1)}_{(k+1)2^{-(n-1)}}\bigr)
  + 2^{-(n+1)/2} Z_{n,k} .
\end{equation}
Join the points of each level linearly. Then with probability one the functions $W^{(n)}$ converge
uniformly on $[0,1]$, and the limit is a standard Brownian motion on $[0,1]$.
\end{theorem}

\begin{proof}
\emph{Finite-dimensional laws.} We show by induction that the increments of $W^{(n)}$ over the $2^n$
intervals of level $n$ are independent $\N(0,2^{-n})$. This holds at level $0$. At level $n$, an old
interval with increment $\Delta\sim\N(0,2^{-(n-1)})$ splits into two increments
$\Delta/2+\sigma Z$ and $\Delta/2-\sigma Z$ with $\sigma^2=2^{-(n+1)}$. Each has variance
$2^{-(n+1)}+2^{-(n+1)}=2^{-n}$, their covariance is $\Var(\Delta)/4-\sigma^2=0$, and being jointly Gaussian
they are independent. Different old intervals use different $Z$'s, which proves the claim. So on the
dyadic grid of level $n$ the values have exactly the laws required by \cref{def3:bm}, and later levels
never change them.

\emph{Uniform convergence.} The function $W^{(n)}-W^{(n-1)}$ is a row of tents. The tent at the
$k$th new point has height $2^{-(n+1)/2}|Z_{n,k}|$, so the sup norm is $2^{-(n+1)/2}\max_k|Z_{n,k}|$. For $c>0$ the Gaussian tail
bound $P(|Z|>x)\le e^{-x^2/2}$, valid for $x\ge1$, gives, for every $n$ with $c\sqrt n\ge1$,
\begin{equation*}
  P\Bigl(\max_{k<2^{n-1}}|Z_{n,k}|>c\sqrt n\Bigr) \le 2^{n-1}e^{-c^2n/2},
\end{equation*}
which is summable in $n$ once $c^2>2\ln 2$; the finitely many levels with $c\sqrt n<1$ do not affect
convergence. By the Borel-Cantelli lemma, with probability one
$\max_k|Z_{n,k}|\le c\sqrt n$ for all large $n$, and then
$\sum_n\lVert W^{(n)}-W^{(n-1)}\rVert_\infty\le\sum_n c\sqrt n\,2^{-(n+1)/2}<\infty$. The partial sums
are continuous and converge uniformly, so the limit $W$ is continuous.

\emph{The limit is Brownian motion.} $W$ agrees with each $W^{(n)}$ on the dyadic points of level $n$,
so it has the right joint laws on all dyadic times. For general times $t_1<\dots<t_k$ take dyadic
$t_j^{(m)}\to t_j$; the vectors $(W_{t^{(m)}_j})_j$ are Gaussian with covariances
$\min(t^{(m)}_i,t^{(m)}_j)$ and converge almost surely by continuity, and a limit of Gaussian vectors
whose covariances converge is Gaussian with the limiting covariance $\min(t_i,t_j)$. A Gaussian process
with mean zero and covariance $\min(s,t)$ has independent increments with the required variances.
\end{proof}

\Cref{fig3:levy} shows one path at levels $0$, $1$, $3$ and $10$. Each level keeps the points of the
previous one and adds detail between them, so a coarse simulation can be refined later without
redrawing it, a property the next section uses. Over \levyPaths{} simulated paths the sample variances
of $W_{1/2}$ and $W_1$ are \levyVarHalf{} and \levyVarOne, the sample value of $\E[W_{1/4}W_{3/4}]$ is
\levyCovQ{} against $\min(\tfrac14,\tfrac34)=0.25$, and the covariance of $W_1-W_{1/2}$ with $W_{1/2}$ is
\levyIncCov.

\begin{figure}[t]
  \centering
  \includegraphics[width=0.86\textwidth]{ch03/fig03_levy}
  \caption{L\'evy's construction. Each level halves the intervals of the previous one and draws the new
  midpoints from the conditional law of \cref{lem3:bridge}; the earlier points never move.}
  \label{fig3:levy}
\end{figure}

\begin{proposition}\label{prop3:bmprops}
Let $W$ be a standard Brownian motion. Then each of the following is also a standard Brownian motion:
$-W$; $c^{-1/2}W_{ct}$ for a constant $c>0$; and $W_{s+t}-W_s$ for a fixed $s\ge 0$, which is also
independent of $(W_r)_{r\le s}$. Each of them, call it $B$, satisfies $\Cov(B_u,B_v)=\min(u,v)$.
\end{proposition}

\begin{proof}
Each transformed process starts at zero and has continuous paths. Each of its increments is plus or
minus a scaled increment of $W$ over a corresponding interval, and disjoint intervals map to disjoint
intervals, so the increments are Gaussian and independent. For scaling,
$\Var\bigl(c^{-1/2}(W_{ct}-W_{cs})\bigr)=c^{-1}c(t-s)=t-s$; the other two variances are immediate. The
independence of $W_{s+\cdot}-W_s$ from the past is property (ii). For the covariance,
$\Cov(B_u,B_v)=\Cov(B_u,B_u+(B_v-B_u))=\Var B_u=u$ for $u\le v$.
\end{proof}

Scaling is the property to remember. Zooming in on a Brownian path by a factor of four in time and two
in space produces another Brownian path, so the path looks equally rough at every scale. That
self-similarity is the intuition behind the next section: the path never smooths out at fine scales,
so it cannot be continuously differentiable on any interval, yet its squared increments add up to
something finite; \cref{sec3:qv} proves both. That a Brownian path is differentiable at no point at
all is a stronger statement with a different proof \parencite[Chapter~1]{morters2010}.

\FloatBarrier
% =====================================================================================
\section{How rough is a path? Quadratic variation}
\label{sec3:qv}

\begin{definition}\label{def3:qv}
For a partition $\Pi=\{0=t_0<t_1<\dots<t_n=T\}$ with mesh $\lVert\Pi\rVert=\max_i(t_{i+1}-t_i)$, the
\emph{quadratic variation} and \emph{total variation} of a path $f$ along $\Pi$ are
\begin{equation}\label{eq3:qvdef}
  Q_\Pi(f)=\sum_{i=0}^{n-1}\bigl(f(t_{i+1})-f(t_i)\bigr)^2,\qquad
  V_\Pi(f)=\sum_{i=0}^{n-1}\bigl|f(t_{i+1})-f(t_i)\bigr| .
\end{equation}
\end{definition}

For a continuously differentiable $f$, $|f(t_{i+1})-f(t_i)|\le\lVert f'\rVert_\infty(t_{i+1}-t_i)$, so, using
$\sum_i(t_{i+1}-t_i)^2\le\lVert\Pi\rVert\,T$,
$Q_\Pi(f)\le\bigl(\lVert f'\rVert_\infty\bigr)^2\,T\,\lVert\Pi\rVert\to0$. Brownian motion behaves differently.

\begin{theorem}\label{thm3:qv}
For Brownian motion and any partition $\Pi$ of $[0,T]$,
\begin{equation}\label{eq3:qvmoments}
  \E\bigl[Q_\Pi(W)\bigr]=T,\qquad
  \E\bigl[(Q_\Pi(W)-T)^2\bigr]=2\sum_{i}(t_{i+1}-t_i)^2\le 2T\lVert\Pi\rVert .
\end{equation}
Hence $Q_\Pi(W)\to T$ in mean square as $\lVert\Pi\rVert\to0$, and almost surely along the dyadic
partitions $t_i=iT2^{-n}$.
\end{theorem}

\begin{proof}
Write $\Delta_i=W_{t_{i+1}}-W_{t_i}$ and $h_i=t_{i+1}-t_i$. Then $\E[\Delta_i^2]=h_i$, which gives the
mean. The terms $Y_i=\Delta_i^2-h_i$ are independent with mean zero, and since $\E[\Delta_i^4]=3h_i^2$,
$\E[Y_i^2]=3h_i^2-2h_i^2+h_i^2=2h_i^2$. So
\begin{equation*}
  \E\Bigl[\Bigl(\sum_i Y_i\Bigr)^2\Bigr]=\sum_i\E[Y_i^2]=2\sum_i h_i^2\le 2\lVert\Pi\rVert\sum_i h_i = 2T\lVert\Pi\rVert .
\end{equation*}
For the dyadic partitions the bound is $2T^2 2^{-n}$. Chebyshev's inequality gives
$P(|Q_n-T|>\varepsilon)\le 2T^22^{-n}/\varepsilon^2$, which is summable, so by Borel-Cantelli
$|Q_n-T|\le\varepsilon$ for all large $n$ almost surely. Taking $\varepsilon=1/k$ for $k=1,2,\dots$ gives
almost sure convergence.
\end{proof}

\begin{corollary}\label{cor3:tv}
Along the dyadic partitions $V_n(W)\to\infty$ almost surely. In particular, with probability one a
Brownian path has infinite total variation on every interval and is continuously differentiable on no
interval.
\end{corollary}

\begin{proof}
On every path $Q_n\le\max_i|\Delta_i|\cdot V_n$. The left side tends to $T>0$, and $\max_i|\Delta_i|\to0$
because a continuous function on $[0,T]$ is uniformly continuous. Hence $V_n\to\infty$. A continuously
differentiable function has finite total variation, which rules it out on $[0,T]$. Applying the
same argument to $W_{s+\cdot}-W_s$ on $[0,q-s]$ for all rationals $s<q$, a countable family, gives
the claim on every interval simultaneously, since every interval contains one with rational ends.
\end{proof}

\Cref{fig3:variation} shows both effects on one simulated path built by \cref{thm3:levy}; the yield
data come at the end of this section. Refining the same path,
the quadratic variation is \qvPathFour{} with $2^4$ intervals, \qvPathEight{} with $2^8$, \qvPathTwelve{}
with $2^{12}$ and \qvPathSixteen{} with $2^{16}$. The total variation over the same partitions is
\tvPathFour, \tvPathEight, \tvPathTwelve{} and \tvPathSixteen; its expected value at $2^{16}$ intervals is
$\sqrt{2n/\pi}=\tvMeanSixteen$, and it doubles every time $n$ is multiplied by four. Panel (b) checks
\cref{eq3:qvmoments}: across simulated paths, the root mean square of $Q_n-1$ follows $\sqrt{2/n}$, with
a fitted slope of \qvSlope{} and no point more than \qvMSEratio{} away from the formula in relative
terms.

\begin{figure}[t]
  \centering
  \includegraphics[width=\textwidth]{ch03/fig04_variation}
  \caption{(a) Quadratic variation (left scale) and total variation (right scale, logarithmic) of one
  Brownian path on $[0,1]$ as the partition is refined. (b) Root mean square error of the quadratic
  variation across simulated paths against $\sqrt{2/n}$.}
  \label{fig3:variation}
\end{figure}

\paragraph{Does the sampling frequency matter for the five year yield?}
If the yield behaved like $\sigma W_t$ with constant $\sigma$, the sum of its squared changes over a
year would estimate $\sigma^2$ without bias at any sampling frequency, daily, weekly or month-end,
with less noise the finer the grid. \Cref{fig3:canqv} tests this on \canN{} daily observations of the five year Government
of Canada yield from \canStart{} to \canEnd. Over the 25 complete calendar years the median realised
volatility, the square root of the year's sum of squared changes, is \canRvDailyMed{} points from
daily changes, \canRvWeeklyMed{} from weekly and \canRvMonthlyMed{} from month-end changes. Year by
year the month-end estimate lies between \canRvRatioLo{} and \canRvRatioHi{} times the daily one in
eight years out of ten. If the daily estimate were exact and the twelve month-end changes independent and normal, the ratio would lie between
\canRvNoiseLo{} and \canRvNoiseHi{} in eight years out of ten, so the observed scatter is about what
sampling noise predicts. The daily estimate
itself ranges from \canRvMin{} points in \canRvMinYear{} to \canRvMax{} in \canRvMaxYear, a wide
spread. Sampling noise contributes to it, but the spread is consistent with volatility that changes
from year to year, though the data alone do not identify how. Panel (b) shows the running sum of squared daily changes through
2025. It ends at \canQvTwentyFive{} squared points and never strays from the straight line to that total
by more than \canQvLinDev\% of the total, and no single day dominates, since the largest squared change, \canQvMaxSq, is about \canQvMaxShare\%
of the year's total: the running sum is close to linear
through 2025. \Cref{thm3:qv} describes the limit for Brownian motion as the grid becomes finer and
says nothing directly about the shape of one year of a yield, so this is a comparison and not a test.

\begin{figure}[t]
  \centering
  \includegraphics[width=\textwidth]{ch03/fig14_canada_qv}
  \caption{Realised quadratic variation of the five year Government of Canada yield. (a) Realised
  volatility (square root of the sum of squared changes) in each calendar year from daily, weekly and
  month-end data. (b) The running sum of squared daily changes in 2025 against a straight line.}
  \label{fig3:canqv}
\end{figure}

The theorem is usually summarised by a multiplication table. It is shorthand for limits of sums over
partitions; no individual increment satisfies it.
\begin{equation}\label{eq3:table}
  (dW_t)^2 = dt,\qquad dW_t\,dt = 0,\qquad (dt)^2 = 0 .
\end{equation}
The first entry says that summed over a fine partition, squared Brownian increments are no longer
random. The other two say that products of higher order vanish faster than $dt$. The $-\tfrac12T$ in
$\int W\,dW$, the $\tfrac12\sigma^2f_{xx}$ in It\^o's lemma and the convexity term of a bond are all
this entry at work.

\subsection{Will the yield touch a level between observation dates? The reflection principle}

The reflection argument of \cref{prop3:srwrefl} has a continuous version. Its proof needs the strong
Markov property: if $\tau$ is a stopping time, then $(W_{\tau+t}-W_\tau)_{t\ge0}$ is again a Brownian
motion independent of the path up to $\tau$. We take that property from \textcite[Section~2.6]{karatzas1998}.

\begin{proposition}\label{prop3:bmrefl}
Let $M_T=\max_{t\le T}W_t$ and $a>0$. Then $P(M_T\ge a)=2P(W_T\ge a)=2\bigl(1-\Phi(a/\sqrt T)\bigr)$, and
the first passage time $\tau_a=\inf\{t:W_t=a\}$ has density
\begin{equation}\label{eq3:fpt}
  f_{\tau_a}(t) = \frac{a}{\sqrt{2\pi t^3}}\exp\Bigl(-\frac{a^2}{2t}\Bigr),\qquad t>0 .
\end{equation}
\end{proposition}

\begin{proof}
Since $W_T\ge a$ implies $M_T\ge a$,
$P(M_T\ge a)=P(W_T\ge a)+P(M_T\ge a,\,W_T<a)$. On $\{\tau_a\le T\}$ the process
$B_t=W_{\tau_a+t}-a$ is a Brownian motion independent of $\tau_a$, and so is $-B$. Reflecting the path
after $\tau_a$ therefore preserves its law, and it maps $\{\tau_a\le T,\,W_T<a\}$ onto
$\{\tau_a\le T,\,W_T>a\}=\{W_T>a\}$. Both events have the same probability, which proves the first claim,
using $P(W_T=a)=0$. Since $\{\tau_a\le t\}=\{M_t\ge a\}$, differentiating
$P(\tau_a\le t)=2\Phi(-a/\sqrt t)$ in $t$ gives
$2\phi(a/\sqrt t)\cdot\tfrac12 a t^{-3/2}$, which is \cref{eq3:fpt}.
\end{proof}

For $a=1$ and $T=1$ the probability is \bmMaxCont\%. A simulation that only looks at the path on a
grid of $n$ dates will miss crossings between grid points and understate it; apart from simulation noise, the gap from
the reflection-principle value measures those missed crossings. That is the right answer
for a contract that checks the barrier only on those dates (discrete monitoring), and the wrong one for
a continuous barrier that is merely observed at those dates. Panel (b) of
\cref{fig3:reflection} shows the bias: with 16 dates the simulated value is \bmMaxMCSixteen\%, and even
with 1{,}024 dates it is \bmMaxMCBig\%, with a standard error of \bmMaxSeBig{} points from 40{,}000 paths
drawn afresh for each grid. \textcite{broadie1997} showed that for monitoring at the $n$ equally spaced dates
$t_k=kT/n$, $k=1,\dots,n$, the discretely monitored probability $P(\max_kW_{t_k}\ge a)$ of this
driftless Brownian motion is, to leading order as $n$ grows, the continuous one with the level moved
up by $\beta\sigma\sqrt{T/n}$ (here $\sigma=1$), where
$\beta=-\zeta(\tfrac12)/\sqrt{2\pi}=\bgkBeta$ and $\zeta$ is the Riemann zeta function. The corrected
formula gives \bmMaxBGKSixteen\% at 16 dates and \bmMaxBGKBig\% at 1{,}024. It is an asymptotic result,
so it is less accurate when there are very few dates. From 128 dates up the gap between the formula and the
simulation is at most \bmMaxZmax{} standard errors, computed from the unrounded values; the largest gap is at
\bmMaxZAt{} dates. The same correction prices discretely monitored barrier options
in Chapter~5.

The five year yield gives a market version of the question. How likely is it to rise half a
percentage point above today's level at some point in the next year? With no drift a fall of the same
size has the same answer, by symmetry. Pooling all daily changes since 2001 gives $\sigma=\canSigma$
points a year. That is higher than the median year's \canRvDailyMed{} of \cref{sec3:qv}, because
pooling gives the high-volatility years greater weight in the sum of squared changes. The
reflection principle then gives $2\bigl(1-\Phi(0.5/\sigma)\bigr)=\canHitModel\%$ for one direction,
twice the \canEndModel\% chance of ending the year that far above. The chance of touching one level
or the other is much higher, \canHitEitherModel\% in the model and \canHitEitherEmp\% in the data. Over
\canStarts{} start dates with a full year of data ahead, the yield went at least half a point up at
some point in \canHitUp\% of the windows and down in \canHitDn\%, an average of \canHitEmp\%. The two
directions differ, plausibly because rates trended over the sample, although a gap of this size is
also within the sampling error described below. Averaging them cancels a steady drift to
first order but cannot remove a trend that changes, and the windows overlap almost completely, so the
effective sample is far smaller than \canStarts{} windows. Its size depends on the serial dependence
of the yield; about \canYearsIndep{} non-overlapping years is only a rough guide, and each
share carries several points of sampling error.

Checking only every 21 trading days, roughly the month-ends, finds a touch in \canHitMonthEmp\% of the
windows on average over the two directions, the discrete monitoring bias of panel (b) in data. The ratio of touching to ending beyond the
level, averaged the same way, is \canHitRatio{} in the data against 2 in the model. Time-varying volatility, drift, mean
reversion and dependence across the overlapping windows could each contribute to the gap, and these
observations do not separate their effects.

\begin{figure}[t]
  \centering
  \includegraphics[width=\textwidth]{ch03/fig02_reflection}
  \caption{(a) A path that reaches the level $a$, and its reflection after the first passage time
  $\tau_a$; both are equally likely. (b) $P(\max_{t\le1}W_t\ge1)$ estimated from paths observed at $n$
  dates (with two standard errors), the continuity corrected formula, and the continuous value.}
  \label{fig3:reflection}
\end{figure}

\FloatBarrier
% =====================================================================================
\section{Trading gains: the It\^o integral}
\label{sec3:ito}

A trading gain is a sum of positions times price changes. If the position $\varphi_t$ is chosen using
information up to $t$ and the price is $W$, the gain over $[0,T]$ should be
$\int_0^T\varphi_t\,dW_t$. Because $W$ has infinite total variation, this cannot be a Stieltjes integral
defined path by path. It\^o's construction uses the probability structure instead, and the left-point
rule that comes with it is the natural one for trading: the position over $[t_i,t_{i+1})$ is fixed at
$t_i$, before the price change is known.

Fix a filtration $(\mathcal F_t)$ to which $W$ is adapted and such that $W_u-W_t$ is independent of
$\mathcal F_t$ for $u>t$. A process $\varphi$ is \emph{adapted} if $\varphi_t$ is
$\mathcal F_t$-measurable for every $t$. Write $L^2_{\mathrm{ad}}$ for the progressively measurable processes (for each $t$ the map
$(s,\omega)\mapsto\varphi_s(\omega)$ on $[0,t]\times\Omega$ is $\mathcal B([0,t])\otimes\mathcal F_t$-measurable) with $\E\int_0^T\varphi_t^2\,dt<\infty$, two of them
being identified when they agree $dt\otimes dP$-almost everywhere. This is the space $L^2(dt\otimes dP)$
restricted to such processes.

\begin{definition}\label{def3:simple}
A process is \emph{simple} if $\varphi_t=\sum_{i=0}^{n-1}\varphi_i\,\mathbf 1_{[t_i,t_{i+1})}(t)$ for a
partition of $[0,T]$ and random variables $\varphi_i$ that are $\mathcal F_{t_i}$-measurable with
$\E[\varphi_i^2]<\infty$. Its It\^o integral is
\begin{equation}\label{eq3:simpleint}
  I(\varphi) = \int_0^T\varphi_t\,dW_t = \sum_{i=0}^{n-1}\varphi_i\,(W_{t_{i+1}}-W_{t_i}) .
\end{equation}
\end{definition}

\begin{proposition}[It\^o isometry for simple processes]\label{prop3:isometry}
For a simple process $\varphi$,
\begin{equation}\label{eq3:isometry}
  \E\bigl[I(\varphi)\bigr]=0,\qquad
  \E\bigl[I(\varphi)^2\bigr]=\E\Bigl[\int_0^T\varphi_t^2\,dt\Bigr],
\end{equation}
and $t\mapsto\int_0^t\varphi_s\,dW_s$ is a martingale.
\end{proposition}

\begin{proof}
Write $\Delta_i=W_{t_{i+1}}-W_{t_i}$ and $h_i=t_{i+1}-t_i$. Since $\varphi_i$ is known at $t_i$ and
$\Delta_i$ is independent of $\mathcal F_{t_i}$,
$\E[\varphi_i\Delta_i]=\E\bigl[\varphi_i\,\E[\Delta_i\mid\mathcal F_{t_i}]\bigr]=0$, which gives the mean.
For the second moment,
\begin{align*}
  \E\bigl[I(\varphi)^2\bigr]
  &= \sum_i\E\bigl[\varphi_i^2\Delta_i^2\bigr] + 2\sum_{i<j}\E\bigl[\varphi_i\varphi_j\Delta_i\Delta_j\bigr] .
\end{align*}
In a cross term with $i<j$, the product $\varphi_i\varphi_j\Delta_i$ is $\mathcal F_{t_j}$-measurable,
so conditioning on $\mathcal F_{t_j}$ gives
$\E[\varphi_i\varphi_j\Delta_i\,\E[\Delta_j\mid\mathcal F_{t_j}]]=0$. In a diagonal term,
$\E[\varphi_i^2\Delta_i^2]=\E\bigl[\varphi_i^2\,\E[\Delta_i^2\mid\mathcal F_{t_i}]\bigr]=\E[\varphi_i^2]h_i$.
Summing, $\E[I^2]=\sum_i\E[\varphi_i^2]h_i=\E\int_0^T\varphi_t^2\,dt$. The martingale property follows
the same way: for $s<t$ every increment after $s$ has conditional mean zero given $\mathcal F_s$, after
conditioning step by step as above.
\end{proof}

\begin{theorem}\label{thm3:itoext}
For every $\varphi\in L^2_{\mathrm{ad}}$, an equivalence class in $L^2(dt\otimes dP)$, there are simple $\varphi^{(n)}$ with
$\E\int_0^T(\varphi_t-\varphi_t^{(n)})^2\,dt\to0$. The integrals $I(\varphi^{(n)})$ converge in mean square
to a limit $I(\varphi)=\int_0^T\varphi_t\,dW_t$ that does not depend on the approximating sequence,
and \cref{eq3:isometry} and the martingale property continue to hold.
\end{theorem}

\begin{proof}[Outline]
The existence of approximations is a density statement about $L^2_{\mathrm{ad}}$, proved in
\textcite[Section~3.2]{karatzas1998}; for adapted $\varphi$ with continuous paths and $\E[\sup_{t\le T}\varphi_t^2]<\infty$
the left-point step functions $\varphi_{t_i}$ on finer and finer partitions will do, by dominated
convergence. For a general element of $L^2_{\mathrm{ad}}$ the density result needs a separate
argument, which the reference gives. Given the approximations, linearity and
\cref{prop3:isometry} give
$\E[(I(\varphi^{(n)})-I(\varphi^{(m)}))^2]=\E\int_0^T(\varphi^{(n)}_t-\varphi^{(m)}_t)^2\,dt$, which tends to
zero. So $(I(\varphi^{(n)}))$ is a Cauchy sequence in mean square and has a limit. Two approximating
sequences can be interleaved into one, so the limit is unique. Means, second moments and conditional
expectations pass to mean square limits, so the isometry and the martingale property survive.
\end{proof}

\begin{proposition}\label{prop3:wdw}
$\displaystyle\int_0^T W_t\,dW_t = \tfrac12\bigl(W_T^2-T\bigr)$.
\end{proposition}

\begin{proof}
For any partition, the identity $b(c-b)=\tfrac12\bigl(c^2-b^2-(c-b)^2\bigr)$ applied with $b=W_{t_i}$
and $c=W_{t_{i+1}}$ and summed over $i$ gives
\begin{equation}\label{eq3:telescope}
  \sum_i W_{t_i}\bigl(W_{t_{i+1}}-W_{t_i}\bigr) = \tfrac12 W_T^2 - \tfrac12 Q_\Pi(W) .
\end{equation}
The left side is the integral of the left-point step approximation of $W$, which converges in mean
square to $\int_0^TW_t\,dW_t$ by \cref{thm3:itoext}. The right side converges in mean square to
$\tfrac12(W_T^2-T)$ by \cref{thm3:qv}.
\end{proof}

Ordinary calculus would give $\tfrac12 W_T^2$. The extra $-\tfrac12T$ is the price of the left-point
rule, and it is what makes the integral a martingale with mean zero. The choice of evaluation point
matters, as the next result shows.

\begin{proposition}\label{prop3:strat}
Let $m_i=(t_i+t_{i+1})/2$. Along partitions with $\lVert\Pi\rVert\to0$,
\begin{equation}\label{eq3:strat}
  \sum_i W_{m_i}\bigl(W_{t_{i+1}}-W_{t_i}\bigr)\ \longrightarrow\ \tfrac12 W_T^2
  \quad\text{in mean square.}
\end{equation}
\end{proposition}

\begin{proof}
Write $A_i=W_{m_i}-W_{t_i}$ and $B_i=W_{t_{i+1}}-W_{m_i}$, independent $\N(0,h_i/2)$. The midpoint sum
equals the left-point sum plus $\sum_i A_i(A_i+B_i)=\sum_iA_i^2+\sum_iA_iB_i$. The first part is a
quadratic variation over half intervals; as in \cref{thm3:qv} it has mean $\sum_ih_i/2=T/2$ and variance
$2\sum_i(h_i/2)^2\to0$. The second has mean zero and variance $\sum_i(h_i/2)^2\to0$. So the extra
terms converge to $T/2$, and \cref{prop3:wdw} completes the proof.
\end{proof}

The limit in \cref{eq3:strat} is the \emph{Stratonovich} integral, written $\int_0^TW_t\circ dW_t$. It
obeys the ordinary chain rule. In this Brownian example it is not a martingale, and midpoint
evaluation uses $W_{m_i}$, which is not yet known at $t_i$ when a trading position must be set. The
left-point rule matches a position set before the price moves. Stratonovich integrals arise instead
as the limit of equations driven by smooth approximations to noise, which is why they suit physical
models.

\Cref{fig3:itostrat}(a) shows both sums on one path with $W_1=\itoWT$, refined from 4 to \itoBigN{}
intervals. With \itoSmallN{} intervals the left-point sum is \itoSumSmall{} and the midpoint sum
\stratSumSmall; with \itoBigN{} they are \itoSumBig{} and \stratSumBig, against limits of \itoLimit{}
and \stratLimit. Across \itoPaths{} paths with 500 intervals, the left-point sums have mean \itoMean{}
and variance \itoVar, close to the isometry value $\int_0^1t\,dt=\tfrac12$, and their histogram matches the
law of $(W_1^2-1)/2$ in panel (b). The midpoint sums exceed them by \stratGapMean{} on average, as
\cref{prop3:strat} predicts. The isometry also holds for other integrands: the sample variance of
$\int_0^1t\,dW_t$ is \isoTVar{} against $\int_0^1t^2dt=\isoTExact$, and the sample second moment of
$\int_0^1e^{W_t}dW_t$ is \isoEVar{} (standard error \isoESE) against $\int_0^1\E[e^{2W_t}]dt=(e^2-1)/2=\isoEExact$.

\begin{figure}[t]
  \centering
  \includegraphics[width=\textwidth]{ch03/fig05_ito_strat}
  \caption{(a) Left-point and midpoint sums for $\int_0^1W\,dW$ on one path as the partition is
  refined; dashed lines mark $(W_1^2-1)/2$ and $W_1^2/2$. (b) Left-point sums across \itoPaths{} paths
  and the density of $(W_1^2-1)/2$.}
  \label{fig3:itostrat}
\end{figure}

\begin{definition}\label{def3:itoprocess}
An \emph{It\^o process} is $X_t=X_0+\int_0^t\mu_s\,ds+\int_0^t\sigma_s\,dW_s$ with adapted $\mu$ and
$\sigma$ such that $\int_0^T|\mu_s|\,ds<\infty$ almost surely and $\sigma\in L^2_{\mathrm{ad}}$. In differential notation,
$dX_t=\mu_t\,dt+\sigma_t\,dW_t$. Its \emph{quadratic variation} is the process $[X]_t=\int_0^t\sigma_s^2\,ds$.
\end{definition}

\begin{remark}
Along partitions of $[0,t]$ with mesh tending to zero, the sums of squared increments of an It\^o
process converge in probability to $[X]_t$. For bounded $\mu$ and $\sigma$ the second moment argument of \cref{thm3:qv}
carries over, with conditioning on $\mathcal F_{t_i}$ in place of independence as in \cref{lem3:qvweighted}, the drift and cross terms vanishing in the limit; the general case follows by stopping
\parencite[Section~1.5]{karatzas1998}. This justifies the name, and nothing below depends on a rate of
convergence.
\end{remark}

\FloatBarrier
% =====================================================================================
\section{It\^o's lemma and what convexity earns}
\label{sec3:lemma}

It\^o's lemma is the chain rule for It\^o processes. Below it gives the convexity term in a bond's expected return, and
Chapter~5 uses it to derive the Black-Scholes equation.

\begin{theorem}[It\^o's lemma]\label{thm3:itolemma}
Let $dX_t=\mu_t\,dt+\sigma_t\,dW_t$ and let $f(t,x)$ have continuous partial derivatives $f_t$, $f_x$
and $f_{xx}$. Then
\begin{equation}\label{eq3:itolemma}
  df(t,X_t) = \Bigl(f_t + \mu_t f_x + \tfrac12\sigma_t^2 f_{xx}\Bigr)dt + \sigma_t f_x\,dW_t ,
\end{equation}
with all derivatives evaluated at $(t,X_t)$.
\end{theorem}

The next lemma controls the accumulated error when squared Brownian increments are replaced by elapsed
time, which is what \cref{eq3:table} abbreviates.

\begin{lemma}\label{lem3:qvweighted}
Let $g$ be adapted with continuous paths and $|g|\le C$. Along partitions with $\lVert\Pi\rVert\to0$,
\begin{equation*}
  \sum_i g_{t_i}\bigl[(W_{t_{i+1}}-W_{t_i})^2-(t_{i+1}-t_i)\bigr]\ \longrightarrow\ 0
  \quad\text{in mean square.}
\end{equation*}
\end{lemma}

\begin{proof}
Let $Y_i=g_{t_i}(\Delta_i^2-h_i)$. For $i<j$, conditioning on $\mathcal F_{t_j}$ shows $\E[Y_iY_j]=0$,
because $\E[\Delta_j^2-h_j\mid\mathcal F_{t_j}]=0$. So
$\E[(\sum_iY_i)^2]=\sum_i\E[g_{t_i}^2]\,\E[(\Delta_i^2-h_i)^2]\le C^2\sum_i2h_i^2\le2C^2T\lVert\Pi\rVert\to0$.
\end{proof}

\begin{proof}[Proof of \cref{thm3:itolemma} in outline]
Take $\mu$ and $\sigma$ constant, so that $\Delta X_i=\mu h_i+\sigma\Delta W_i$ exactly on every interval. Fix
$K$ and let $f_K$ agree with $f$ on $[0,t]\times[-K,K]$ and vanish for $\lvert x\rvert\ge K+1$ (multiply $f$ by
a smooth cutoff in $x$). Then $f_{K,t}$, $f_{K,x}$ and $f_{K,xx}$ are continuous with compact support, so
they are bounded and uniformly continuous. We prove the identity for $f_K$, where
\cref{lem3:qvweighted} applies to $g_{t_i}=f_{K,xx}(t_i,X_{t_i})$, and write $f$ for $f_K$ until the last
step. On a partition of $[0,t]$ write $\Delta X_i=X_{t_{i+1}}-X_{t_i}$.
Taylor's theorem, applied first in $t$ and then in $x$, gives
\begin{equation*}
  f(t,X_t)-f(0,X_0) = \sum_i\Bigl[f_t\,h_i + f_x\,\Delta X_i + \tfrac12 f_{xx}\,(\Delta X_i)^2\Bigr] + R_\Pi ,
\end{equation*}
with derivatives at $(t_i,X_{t_i})$ and a remainder $R_\Pi$ bounded at the end of the argument.
Now take limits term by term. The first sum is a Riemann sum for $\int_0^tf_s\,ds$. The second is
$\sum f_x\mu h_i+\sum f_x\sigma\Delta W_i$, which tends to $\int f_x\mu\,ds+\int f_x\sigma\,dW_s$, the
first as a Riemann sum and the second by \cref{thm3:itoext}. In the third,
$(\Delta X_i)^2=\sigma^2\Delta W_i^2+2\mu\sigma h_i\Delta W_i+\mu^2h_i^2$; the last two pieces sum to
something that tends to zero in mean square, and by \cref{lem3:qvweighted}
$\sum\tfrac12f_{xx}\sigma^2\Delta W_i^2$ has the same limit as $\sum\tfrac12f_{xx}\sigma^2h_i$, namely
$\int_0^t\tfrac12\sigma^2f_{xx}\,ds$. Finally $|R_\Pi|\le\varepsilon_\Pi\sum_i\bigl(h_i+(\Delta X_i)^2\bigr)$, where $\varepsilon_\Pi$ is the largest
modulus of continuity of $f_t$ and $f_{xx}$ over the increments, which tends to zero almost surely by
uniform continuity, since the largest increment of the continuous path tends to zero. The sum converges to $t+\sigma^2t$ in
probability, so $R_\Pi\to0$ in probability. Every other term converges in mean square, hence in
probability, so for each $t$ the identity holds with probability one, by passing to a subsequence
along which all the convergences are almost sure. Both sides are continuous in $t$, so it holds for
all rational $t$ at once and then for every $t$. Now return to $f$. On the event $\{\max_{s\le t}\lvert X_s\rvert< K\}$ the time and drift integrals
for $f_K$ and for $f$ coincide, and so do the stochastic integrals, because their integrands agree
up to the time $X$ first reaches $\pm K$ and an integral stopped at that time depends only on the
integrand before it. The path is continuous, so this event has probability tending to one as
$K\to\infty$, and letting $K\to\infty$ removes the localisation. For
general adapted $\mu$ and $\sigma$ the drift sum is no longer a Riemann sum and has to be approximated in
its own integrability sense, and the stochastic integral needs the $L^2$ approximation by step
integrands after localisation; we do not carry this out \parencite[Section~3.3]{karatzas1998}.
\end{proof}

The extra term $\tfrac12\sigma^2f_{xx}\,dt$ is the whole difference from ordinary calculus. The first
use is a bond.

\begin{example}{What convexity earns a five year bond}\label{ex3:convexity}
Let a five year zero coupon bond trade at $P=e^{-5y}$, with $y$ its continuously compounded yield
in decimal, and suppose over a short horizon $dy=\sigma_y\,dW$. It\^o's lemma with $f(y)=e^{-5y}$ gives
$dP/P=-5\,dy+\tfrac12\cdot25\,\sigma_y^2\,dt$. A holder earns the second term on average even when
the yield does not move on average, because a fall in yield raises the price by more than an equal rise
lowers it. Take the annualised volatility of the five year Government of Canada yield's daily changes since
2001, $\canSigma$ points, so that $\sigma_y=0.0082$ in decimal. The convexity term is then
$12.5\,\sigma_y^2=\canConvex\%$ a year, about \canConvexBp{} basis points. The benchmark is a
semi-annually compounded par yield, so this is a model-implied figure for an idealised zero coupon
bond of the same volatility, an order of magnitude for the convexity contribution and no promise of
what a holder earns. Chapter~4 returns to convexity in the term structure.
\end{example}

The next three examples apply the formula to familiar functions of Brownian motion alone, to show what its second derivative term does before the lemma meets more general models.

\begin{example}{Powers of Brownian motion}\label{ex3:powers}
With $f(x)=x^3$ and $X=W$, \cref{eq3:itolemma} gives $d(W_t^3)=3W_t^2\,dW_t+3W_t\,dt$, that is
$W_T^3=\int_0^T3W_t^2\,dW_t+3\int_0^TW_t\,dt$. Taking expectations, $\E[W_T^3]=0$, as symmetry requires.
With $f(x)=x^2$, $d(W_t^2)=2W_t\,dW_t+dt$, which is \cref{prop3:wdw} again, and shows that $W_t^2-t$ is a
martingale.
\end{example}

\begin{example}{The exponential martingale}\label{ex3:expmart}
Let $Z_t=\exp(\theta W_t-\tfrac12\theta^2t)$ and $f(t,x)=\exp(\theta x-\tfrac12\theta^2t)$. Then
$f_t=-\tfrac12\theta^2f$, $f_x=\theta f$ and $f_{xx}=\theta^2f$, so the $dt$ terms cancel and
$dZ_t=\theta Z_t\,dW_t$. Since $\E[Z_t^2]=e^{\theta^2t}$ is finite, the integrand $\theta Z$ lies in
$L^2_{\mathrm{ad}}$ on every $[0,T]$, so $Z$ is a true martingale; it is the stochastic exponential that
\cref{sec3:girsanov} uses to change measure.
\end{example}

\begin{example}{Integration by parts}\label{ex3:parts}
For $f(t,x)=tx$, $f_{xx}=0$ and $d(tW_t)=W_t\,dt+t\,dW_t$. Hence
$\int_0^TW_t\,dt=TW_T-\int_0^Tt\,dW_t=\int_0^T(T-t)\,dW_t$, a Gaussian variable whose variance, by the
isometry, is $\int_0^T(T-t)^2dt=T^3/3$.
\end{example}

The two-dimensional version is proved the same way. Let $dX=\mu^X_t\,dt+\sigma^X_t\,dW_t$ and
$dY=\mu^Y_t\,dt+\sigma^Y_t\,dB_t$, where $W$ and $B$ are Brownian motions with instantaneous correlation
$\rho$, that is $[W,B]_t=\rho t$. Then
\begin{equation}\label{eq3:product}
  d(X_tY_t) = X_t\,dY_t+Y_t\,dX_t+d[X,Y]_t,\qquad d[X,Y]_t=\sigma^X_t\sigma^Y_t\rho\,dt ,
\end{equation}
and more generally $df(X,Y)$ has the second-order terms
$\tfrac12f_{xx}d[X]+f_{xy}d[X,Y]+\tfrac12f_{yy}d[Y]$.

\Cref{fig3:itolemma}(a) tests \cref{eq3:itolemma} on simulated paths. For $f(w)=w^3$ and
$f(w)=e^w$ it compares $f(W_1)-f(0)$ with the discretised right side, with and without the second-order
term. With the term included the root mean square residual falls like $n^{-1/2}$ (fitted slopes
\lemmaSlopeCube{} and \lemmaSlopeExp) and is \lemmaRmsCube{} and \lemmaRmsExp{} at \lemmaN{}
intervals. Without it, the residual stays near \naiveRmsCube{} and \naiveRmsExp{} however fine the grid.
Panel (b) shows what the ordinary chain rule misses for $w^3$: in the limit it is exactly
$3\int_0^1W_t\,dt$, path by path, and on the 4,096-interval grid the two have correlation \naiveGapCorr.

\begin{figure}[t]
  \centering
  \includegraphics[width=\textwidth]{ch03/fig06_ito_lemma}
  \caption{(a) Root mean square residual of It\^o's lemma and of the ordinary chain rule for $W^3$ and $e^W$,
  over \lemmaPaths{} paths. (b) For $W^3$, the gap left by the chain rule against $3\int_0^1W_t\,dt$.}
  \label{fig3:itolemma}
\end{figure}

\FloatBarrier
% =====================================================================================
\section{Which level comes first? Martingales and stopping times}
\label{sec3:mart}

\begin{definition}\label{def3:mart}
An adapted process $M$ with $\E|M_t|<\infty$ is a \emph{martingale} if $\E[M_t\mid\mathcal F_s]=M_s$ for
all $s\le t$. A random time $\tau$ is a \emph{stopping time} if $\{\tau\le t\}\in\mathcal F_t$ for every
$t$: whether it has happened by $t$ is known at $t$.
\end{definition}

\begin{proposition}\label{prop3:martingales}
For a Brownian motion $W$ and a constant $\theta$, the processes $W_t$, $W_t^2-t$ and
$Z_t=\exp(\theta W_t-\tfrac12\theta^2t)$ are martingales.
\end{proposition}

\begin{proof}
Write $W_t=W_s+D$ with $D=W_t-W_s\sim\N(0,t-s)$ independent of $\mathcal F_s$. Then
$\E[W_t\mid\mathcal F_s]=W_s$; $\E[W_t^2\mid\mathcal F_s]=W_s^2+2W_s\E[D]+\E[D^2]=W_s^2+t-s$; and
$\E[Z_t\mid\mathcal F_s]=Z_s\,\E[e^{\theta D}]e^{-\theta^2(t-s)/2}=Z_s$, using
$\E[e^{\theta D}]=e^{\theta^2(t-s)/2}$.
\end{proof}

A martingale has constant mean, which is easy to check by simulation and easy to get wrong. Over
\mgPaths{} draws of $W_1$, the sample means of $W_1^2-1$ and of $e^{W_1-1/2}$ are \mgSq{} and \mgExp, as
they should be. For $\theta=3$ the sample mean of $Z_1$ is \mgExpThree{} with a sample standard error of
\mgExpThreeSE. That looks like a clean confirmation, but the standard error is the weak point. The true
standard deviation of $Z_1$ is $\sqrt{e^9-1}=\mgExpThreeSD$, so the true standard error is
\mgExpThreeTrueSE, \mgSERatio{} times the estimate; only \mgExpThreeShare\% of the draws exceed one, and
the mean is held up by rare huge values that a sample of this size mostly misses. Repeating the
experiment \mgReps{} times, each with a fresh sample of the same size, the sample standard error was
below half the true one in \mgSEHalfShare\% of the samples, and a test at two sample standard errors
rejected the true mean of one in \mgRejShare\% of them, against the 5\% it is meant to allow. Heavy
right tails make the sample standard error itself unreliable, a warning that applies to any
simulation of likelihood ratios or of leveraged payoffs.

Stopping a martingale at a stopping time leaves its mean unchanged only under conditions, and without
them it can fail. The optional stopping theorem gives two sufficient sets: if
$M$ is a martingale with continuous paths and $\tau$ is a stopping time, then $\E[M_\tau]=\E[M_0]$ when
$\tau$ is bounded, and also when $\tau<\infty$ almost surely and the stopped process $M_{t\wedge\tau}$
is bounded by a constant \parencite[Section~1.3]{karatzas1998}. A doubling strategy that quits at its first
gain beats a fair game only with unlimited credit, and no lender offers that.

\begin{proposition}\label{prop3:exit}
Let $X_t=\mu t+\sigma W_t$ with $\sigma>0$, let $a,b>0$ and let $\tau$ be the first time $X$ hits $a$ or
$-b$. Then $\tau<\infty$ almost surely, $\E[\tau]<\infty$, and with $k=2\mu/\sigma^2$,
\begin{equation}\label{eq3:exit}
  P(X_\tau=a)=p=\frac{1-e^{kb}}{e^{-ka}-e^{kb}},\qquad
  \E[\tau]=\frac{ap-b(1-p)}{\mu}\quad(\mu\neq0),
\end{equation}
with the limits $p=b/(a+b)$ and $\E[\tau]=ab/\sigma^2$ when $\mu=0$.
\end{proposition}

\begin{proof}
\emph{The exit time is finite.} On each interval $[j,j+1]$ the increment of $X$ exceeds $a+b$ with a
probability $q>0$ that does not depend on $j$, and if that happens the process cannot have stayed
inside $(-b,a)$. The increments over disjoint intervals are independent, so $P(\tau>n)\le(1-q)^n$, and
$\E[\tau]=\int_0^\infty P(\tau>t)\,dt<\infty$.

\emph{The exit probability.} By It\^o's lemma, $Y_t=e^{-kX_t}$ satisfies
$dY_t=Y_t\bigl(-k\mu+\tfrac12k^2\sigma^2\bigr)dt-k\sigma Y_t\,dW_t=-k\sigma Y_t\,dW_t$, because
$\tfrac12k^2\sigma^2=k\mu$. So $Y$ is a martingale, and $Y_{t\wedge\tau}$ lies between $e^{-ka}$ and $e^{kb}$.
Optional stopping gives $1=Y_0=\E[Y_\tau]=pe^{-ka}+(1-p)e^{kb}$, which solves to \cref{eq3:exit}.

\emph{The expected exit time.} $X_t-\mu t=\sigma W_t$ is a martingale. At $t\wedge\tau$,
$\E[X_{t\wedge\tau}]=\mu\E[t\wedge\tau]$. As $t\to\infty$ the left side converges to $\E[X_\tau]=ap-b(1-p)$
by bounded convergence, and the right side to $\mu\E[\tau]$ by monotone convergence. For $\mu=0$ use
the martingales $X_t$ and $X_t^2-\sigma^2t$ in the same way: $0=ap-b(1-p)$ gives $p=b/(a+b)$, and
$\E[X_\tau^2]=\sigma^2\E[\tau]$ gives $a^2p+b^2(1-p)=ab=\sigma^2\E[\tau]$.
\end{proof}

Read $X$ as the log return of a fund since inception, $a$ as a target and $-b$ as a fixed cumulative
loss from inception at which a lender or an investor acts. That is a loss from the starting value,
which differs from a drawdown measured from the running peak. With $a=b=10\%$, $\sigma=20\%$ and $\mu=10\%$, the target is reached first
with probability \exitPup\%, and the race lasts \exitTau{} years on average; with no drift the
probability is one half and the expected time is \exitTauZero{} years. \Cref{fig3:exit} compares the
formulas with \exitPaths{} simulated paths at each drift, in steps of $0.001$ years with a Brownian bridge test for
crossings between steps. At $\mu=10\%$ the simulation gives \exitPupMC\% and \exitTauMC{} years.

\begin{figure}[t]
  \centering
  \includegraphics[width=\textwidth]{ch03/fig12_exit}
  \caption{Exit of $\mu t+\sigma W_t$ from $(-0.1,0.1)$ with $\sigma=0.2$. (a) Probability of reaching the
  upper level first. (b) Expected exit time. Lines are \cref{eq3:exit}; points are simulations.}
  \label{fig3:exit}
\end{figure}

\FloatBarrier
% =====================================================================================
\section{Simulating prices and rates: stochastic differential equations}
\label{sec3:sde}

A stochastic differential equation $dX_t=a(t,X_t)\,dt+b(t,X_t)\,dW_t$ with $X_0=x_0$ is shorthand for
the integral equation $X_t=x_0+\int_0^ta(s,X_s)\,ds+\int_0^tb(s,X_s)\,dW_s$. A \emph{strong solution} is
an adapted continuous process that satisfies it on the given probability space. If $a$ and $b$ are
jointly measurable, Lipschitz in $x$ uniformly in $t$, and grow at most linearly, and $x_0$ is a
constant, there is exactly one strong solution with $\E\int_0^T X_t^2\,dt<\infty$; the proof is a Picard iteration in mean square
\parencite[Theorem~5.2.1]{oksendal2003}. Weaker conditions also suffice; these are enough for this chapter. The two equations below satisfy these conditions and can be
solved in closed form.

\subsection{Geometric Brownian motion}

\begin{proposition}\label{prop3:gbm}
The equation $dS_t=\mu S_t\,dt+\sigma S_t\,dW_t$ with $S_0>0$ has the unique solution
\begin{equation}\label{eq3:gbm}
  S_t = S_0\exp\Bigl(\bigl(\mu-\tfrac12\sigma^2\bigr)t+\sigma W_t\Bigr).
\end{equation}
Hence $\ln S_t\sim\N\bigl(\ln S_0+(\mu-\tfrac12\sigma^2)t,\ \sigma^2t\bigr)$, $\E[S_t]=S_0e^{\mu t}$,
the median of $S_t$ is $S_0e^{(\mu-\sigma^2/2)t}$ and $P(S_t<S_0)=\Phi\bigl(-(\mu-\tfrac12\sigma^2)\sqrt t/\sigma\bigr)$.
\end{proposition}

\begin{proof}
Let $S_t$ be defined by \cref{eq3:gbm}, so $S_t=f(t,W_t)$ with
$f(t,w)=S_0\exp((\mu-\tfrac12\sigma^2)t+\sigma w)$. Then $f_t=(\mu-\tfrac12\sigma^2)f$, $f_w=\sigma f$ and
$f_{ww}=\sigma^2f$, and It\^o's lemma gives
\begin{equation*}
  dS_t = \Bigl(\mu-\tfrac12\sigma^2+\tfrac12\sigma^2\Bigr)S_t\,dt+\sigma S_t\,dW_t = \mu S_t\,dt+\sigma S_t\,dW_t .
\end{equation*}
The coefficients $\mu x$ and $\sigma x$ are Lipschitz, so this is the only solution. The distribution of
$\ln S_t$ is read off \cref{eq3:gbm}, and the mean, median and loss probability follow from it and from
the lognormal moments derived in Chapter~1.
\end{proof}

The drift of the logarithm is $\mu-\tfrac12\sigma^2$, not $\mu$. That correction is the gap between the
arithmetic and geometric mean returns met in Chapter~1. The median of $S_T$ grows at the rate $\mu-\tfrac12\sigma^2$
while its mean grows at $\mu$, so over long horizons the typical path falls far below the expected
value.

\begin{proposition}\label{prop3:gbmlong}
Let $\sigma>0$. If $\mu<\tfrac12\sigma^2$, then $S_t\to0$ almost surely. If in addition $\mu>0$, then at the same
time $\E[S_t]=S_0e^{\mu t}\to\infty$: almost every path vanishes while the average over paths grows.
\end{proposition}

\begin{proof}
By \cref{eq3:gbm}, $\ln S_t=\ln S_0+t\bigl(\mu-\tfrac12\sigma^2+\sigma W_t/t\bigr)$, so it is enough to show
$W_t/t\to0$ almost surely. On integer times $W_n=\sum_{k\le n}(W_k-W_{k-1})$ is a sum of independent
standard normals, so $W_n/n\to0$ by the strong law of large numbers. Between integers, by
\cref{prop3:bmrefl} applied to $W$ and $-W$,
$P\bigl(\max_{n\le t\le n+1}|W_t-W_n|\ge n^{2/3}\bigr)\le4\bigl(1-\Phi(n^{2/3})\bigr)$, which is summable. By
Borel-Cantelli the oscillation over $[n,n+1]$ is below $n^{2/3}$ for all large $n$, and dividing by
$n$ gives the claim.
\end{proof}

\Cref{fig3:gbm} shows the two regimes over thirty years with $S_0=100$ and $\mu=8\%$. With
$\sigma=25\%$ the log drift is \gbmGrowthLow\% a year and the median path grows to \gbmMedianLow. With
$\sigma=50\%$ the log drift is \gbmGrowthHigh\% and the median falls to \gbmMedianHigh, while the mean is
\gbmMeanHigh{} in both cases. In the second case the probability of ending below the start is
\gbmProbLossHigh\%. The mean is carried by a few paths that do extraordinarily well, which is why a
volatile asset with an attractive expected return can still be a poor thing to hold alone, and why the
geometric mean is the better guide to what a single investor will see.

\begin{figure}[t]
  \centering
  \includegraphics[width=\textwidth]{ch03/fig07_gbm}
  \caption{Forty paths of geometric Brownian motion over thirty years, with the mean, the median and the
  central 90\% band. (a) $\sigma=25\%$: the median grows. (b) $\sigma=50\%$: the mean is unchanged but
  the median falls, since $\mu<\sigma^2/2$.}
  \label{fig3:gbm}
\end{figure}

\subsection{Two numerical schemes}

Most stochastic differential equations have no closed form, and the standard way to simulate them is
to step forward on a grid $t_k=kh$. For one Brownian driver, the \emph{Euler-Maruyama} scheme and the scalar-noise \emph{Milstein} scheme are
\begin{align}
  \hat X_{k+1} &= \hat X_k + a(\hat X_k)h + b(\hat X_k)\Delta W_k, \label{eq3:euler}\\
  \hat X_{k+1} &= \hat X_k + a(\hat X_k)h + b(\hat X_k)\Delta W_k + \tfrac12 b(\hat X_k)b'(\hat X_k)\bigl(\Delta W_k^2-h\bigr), \label{eq3:milstein}
\end{align}
with $\Delta W_k=W_{t_{k+1}}-W_{t_k}$. The Milstein term comes from expanding $b(X_t)$ inside the
stochastic integral by It\^o's lemma and keeping the piece of order $h$; by \cref{eq3:table} its
average is zero but its size is not. Two notions of accuracy matter. The \emph{strong} error
$\E|X_T-\hat X_T|$ measures how close simulated paths are to true paths driven by the same noise, which
matters for hedging simulations. The \emph{weak} error $|\E g(X_T)-\E g(\hat X_T)|$ measures how close
expectations are, which is what pricing needs. If $a$ and $b$ are globally Lipschitz with linear growth
and the starting value has a finite second moment, the Euler scheme has strong order $\tfrac12$. The
Milstein scheme has strong order $1$ if, in addition, $a$ and $b$ are smooth enough for the terms below to
exist, $bb'$ is globally Lipschitz, and $bb'$ and the coefficients of the next terms of the It\^o-Taylor
expansion, such as $ab'+\tfrac12b^2b''$ and $b(bb')'$, grow at most linearly. A bounded $b'$ alone is not
enough. The Euler scheme has weak order $1$ if, besides the Lipschitz and growth conditions, $a$, $b$
and the payoff $g$ are four times continuously differentiable with derivatives of polynomial growth
\parencite{kloeden1992,higham2001}. Geometric Brownian motion, with $a(x)=\mu x$ and $b(x)=\sigma x$,
meets all of these conditions except the smoothness of a payoff with a kink, such as a call.

For geometric Brownian motion the Euler mean can be computed exactly, which gives a clean check.

\begin{proposition}\label{prop3:eulerweak}
For the Euler scheme applied to $dS=\mu S\,dt+\sigma S\,dW$ with $n$ steps of size $h=T/n$,
\begin{equation}\label{eq3:eulermean}
  \E[\hat S_T]=S_0(1+\mu h)^n,\qquad
  \E[\hat S_T^2]=S_0^2\bigl((1+\mu h)^2+\sigma^2h\bigr)^n,
\end{equation}
and $\E[\hat S_T]-S_0e^{\mu T}=-\tfrac12S_0e^{\mu T}\mu^2Th+O(h^2)$.
\end{proposition}

\begin{proof}
$\hat S_T=S_0\prod_k(1+\mu h+\sigma\Delta W_k)$ is a product of independent factors with mean
$1+\mu h$ and second moment $(1+\mu h)^2+\sigma^2h$, which gives \cref{eq3:eulermean}. For the expansion,
$n\ln(1+\mu h)=\frac Th\bigl(\mu h-\tfrac12\mu^2h^2+O(h^3)\bigr)=\mu T-\tfrac12\mu^2Th+O(h^2)$, and
exponentiating gives $(1+\mu h)^n=e^{\mu T}\bigl(1-\tfrac12\mu^2Th+O(h^2)\bigr)$.
\end{proof}

With $S_0=100$, $\mu=8\%$ and twelve monthly steps over a year, the Euler mean is \weakMeanTwelve{}
against the exact \exactMeanOne. The leading term of the expansion, $-\tfrac12S_0e^{\mu T}\mu^2Th$, is \weakLead, close to the actual gap.
\Cref{fig3:schemes} shows both orders with $\sigma=40\%$. Panel (a) uses \strongPaths{} paths on a
common fine grid; the fitted strong orders are \strongEuler{} for Euler and \strongMilstein{} for
Milstein, and with 16 steps the errors are \strongEulerSixteen{} and \strongMilSixteen. Panel (b) plots the
exact weak errors of \cref{eq3:eulermean}, with fitted slopes \weakSlopeMean{} and \weakSlopeSecond.

\begin{figure}[t]
  \centering
  \includegraphics[width=\textwidth]{ch03/fig08_schemes}
  \caption{Accuracy of discretisation schemes for geometric Brownian motion with $\mu=8\%$ and
  $\sigma=40\%$ over one year. (a) Strong error against step size. (b) Exact weak error of the Euler
  scheme for the first and second moments.}
  \label{fig3:schemes}
\end{figure}

\subsection{The Ornstein-Uhlenbeck process}

Interest rates, credit spreads and the basis between related commodities are frequently modelled as
mean reverting, although they can shift and trend for years, as Canadian yields did from 2001 to
2020. The simplest mean-reverting model is due to
\textcite{uhlenbeck1930}, and \textcite{vasicek1977} used it for the short rate:
\begin{equation}\label{eq3:ou}
  dX_t = \kappa(\theta-X_t)\,dt+\sigma\,dW_t,\qquad \kappa>0 .
\end{equation}
The drift points towards $\theta$ and is stronger the further away $X$ is.

\begin{proposition}\label{prop3:ou}
The solution of \cref{eq3:ou} with $X_0=x_0$ is
\begin{equation}\label{eq3:ousol}
  X_t = \theta+(x_0-\theta)e^{-\kappa t}+\sigma\int_0^te^{-\kappa(t-s)}\,dW_s .
\end{equation}
It is Gaussian with
\begin{gather}\label{eq3:oumoments}
  \E[X_t]=\theta+(x_0-\theta)e^{-\kappa t},\qquad
  \Var[X_t]=\frac{\sigma^2}{2\kappa}\bigl(1-e^{-2\kappa t}\bigr),\\
  \Cov(X_s,X_t)=\frac{\sigma^2}{2\kappa}\bigl(e^{-\kappa(t-s)}-e^{-\kappa(t+s)}\bigr)\notag
\end{gather}
for $s\le t$. As $t\to\infty$ the law of $X_t$ tends to $\N(\theta,\sigma^2/(2\kappa))$, and the gap between
the mean and $\theta$ halves every $\ln 2/\kappa$ units of time.
\end{proposition}

\begin{proof}
Apply the product rule \cref{eq3:product} to $Y_t=e^{\kappa t}(X_t-\theta)$. The factor $e^{\kappa t}$
has no Brownian part, so there is no cross term, and
\begin{equation*}
  dY_t=\kappa e^{\kappa t}(X_t-\theta)\,dt+e^{\kappa t}\,dX_t
  =\kappa e^{\kappa t}(X_t-\theta)\,dt+e^{\kappa t}\bigl(\kappa(\theta-X_t)\,dt+\sigma\,dW_t\bigr)
  =\sigma e^{\kappa t}\,dW_t .
\end{equation*}
Integrating, $Y_t=x_0-\theta+\sigma\int_0^te^{\kappa s}\,dW_s$, and multiplying by $e^{-\kappa t}$ gives
\cref{eq3:ousol}. The stochastic integral has a deterministic integrand, so it is a mean square limit
of Gaussian sums and is Gaussian with mean zero. By the isometry its variance is
$\sigma^2\int_0^te^{-2\kappa(t-s)}ds=\sigma^2(1-e^{-2\kappa t})/(2\kappa)$. For the covariance, split the
integral up to $t$ at $s$; the part over $(s,t]$ is independent of $X_s$, and the part over $[0,s]$
contributes $\sigma^2e^{-\kappa(t-s)}\int_0^se^{-2\kappa(s-u)}du$, which simplifies to the stated form.
The limits follow from $e^{-\kappa t}\to0$, and $\E[X_t]-\theta=(x_0-\theta)e^{-\kappa t}$ halves when
$e^{-\kappa t}=\tfrac12$.
\end{proof}

\begin{corollary}\label{cor3:ar1}
Observed at spacing $\delta$, the process is exactly a first order autoregression:
\begin{equation}\label{eq3:ar1}
  X_{k+1}=\theta(1-\phi)+\phi X_k+\varepsilon_{k+1},\qquad
  \phi=e^{-\kappa\delta},\qquad \varepsilon_{k+1}\sim\N\Bigl(0,\frac{\sigma^2(1-\phi^2)}{2\kappa}\Bigr)\ \text{i.i.d.}
\end{equation}
\end{corollary}

\begin{proof}
Apply \cref{prop3:ou} over one interval starting from $X_k$, using the independence of the
Brownian increments after time $k\delta$.
\end{proof}

The corollary gives an exact simulation method with no discretisation error, and it is the link to
the smoothing models of Chapter~1, where appraisal-based returns were modelled as autoregressions. It
also gives an estimator: regress $X_{k+1}$ on $X_k$ and set $\hat\kappa=-\ln\hat\phi/\delta$. The listing
below does this on fifty years of simulated monthly data with $\kappa=1.2$.
\lstinputlisting[language=Python]{code/ch03/snippets/ou_fit.py}
It prints $\hat\kappa=\snipKappa$ and a half-life of \snipHalf{} years, against true values of $1.2$ and
\ouHalfLife. The estimate is \snipExcess{} too high. About \ouKendallKappa{} of that excess is the small-sample bias described next; the remaining
\snipExcessRest{} is small beside the sampling noise, and another sample could miss in either direction. \textcite{kendall1954} showed (with the mean
estimated) that the least squares estimate of an autoregressive coefficient is biased towards zero by
about $(1+3\phi)/n$ in a sample of $n$, a first-order approximation for a Gaussian first order
autoregression. With $\phi=\ouPhiMonthly$ and $n=600$ that is \ouKendall, and the
average $\hat\phi$ over 400 simulated samples is off by \ouPhiBias. Converted to $\kappa$, that bias is about \ouKendallKappa{} a year, a
fraction of the sampling standard deviation of $\hat\kappa$ at this length (\ouKappaSd{} in the same
simulations). The map $\hat\kappa=-\ln\hat\phi/\delta$ multiplies relative error by $1/(-\ln\phi)$, here about
\ouAmp, so a \ouPhiBiasRel\% bias in $\hat\phi$ becomes about \ouKappaBiasRel\% in $\hat\kappa$; near $\phi=1$,
as in the case study, the factor is about \caseAmp.

\Cref{fig3:ou} shows thirty paths that start at 10\% with $\kappa=1.2$, $\theta=3\%$ and $\sigma=2\%$. The
mean is \ouMeanOne\% after one year, and the half-life is \ouHalfLife{} years. The stationary standard
deviation is \ouSdInf{} points, so the model gives a negative value \ouProbNeg\% of the time in the long
run. Until 2012 that was counted against the model; once Denmark, then the euro area, Switzerland and
Sweden set negative policy rates, it looked more like a feature. That the Gaussian model allows
negative values does not show that it describes how rates behave near or below zero.

\begin{figure}[t]
  \centering
  \includegraphics[width=0.86\textwidth]{ch03/fig09_ou}
  \caption{Thirty Ornstein-Uhlenbeck paths simulated exactly by \cref{cor3:ar1} on a daily grid, with
  the mean and a 95\% band from \cref{eq3:oumoments}.}
  \label{fig3:ou}
\end{figure}

\FloatBarrier
% =====================================================================================
\section{Changing the odds: Girsanov's theorem}
\label{sec3:girsanov}

Two probability measures $P$ and $Q$ on $\mathcal F_T$ are \emph{equivalent} if they agree on which
events are impossible. Then there is a positive random variable $Z_T=dQ/dP$, the Radon-Nikodym
density, with $\E_Q[Y]=\E_P[YZ_T]$ for every integrable $Y$. Changing measure reweights the paths
without changing which paths can occur. The question is what happens to a Brownian motion under the
new weights.

\begin{theorem}[Girsanov, constant drift]\label{thm3:girsanov}
Let $W$ be a $P$-Brownian motion on $[0,T]$, let $\theta$ be a constant and define $Q$ by
\begin{equation}\label{eq3:girsanov}
  \frac{dQ}{dP}=Z_T=\exp\Bigl(-\theta W_T-\tfrac12\theta^2T\Bigr).
\end{equation}
Then $Q$ is a probability measure equivalent to $P$, and $\tilde W_t=W_t+\theta t$ is a $Q$-Brownian
motion on $[0,T]$.
\end{theorem}

\begin{proof}
$Z_T>0$ and $\E_P[Z_T]=1$ by \cref{prop3:martingales} with $-\theta$ in place of $\theta$, so $Q$ is a
probability measure equivalent to $P$. The paths of $\tilde W$ are continuous and start at zero, so it
remains to find the law of its increments under $Q$. Take $0=t_0<\dots<t_n=T$, write
$\Delta_j=W_{t_j}-W_{t_{j-1}}$, $h_j=t_j-t_{j-1}$, and let $g$ be bounded. Since $W_T=\sum_j\Delta_j$,
the density factorises, $Z_T=\prod_j\exp(-\theta\Delta_j-\tfrac12\theta^2h_j)$, and the $\Delta_j$ are
independent under $P$. Hence
\begin{align*}
  \E_Q\bigl[g(\Delta_1+\theta h_1,\dots,\Delta_n+\theta h_n)\bigr]
  &= \int g(x_1+\theta h_1,\dots,x_n+\theta h_n)\prod_j\frac{e^{-x_j^2/(2h_j)-\theta x_j-\theta^2h_j/2}}{\sqrt{2\pi h_j}}\,dx\\
  &= \int g(x_1+\theta h_1,\dots,x_n+\theta h_n)\prod_j\frac{e^{-(x_j+\theta h_j)^2/(2h_j)}}{\sqrt{2\pi h_j}}\,dx\\
  &= \int g(y_1,\dots,y_n)\prod_j\frac{e^{-y_j^2/(2h_j)}}{\sqrt{2\pi h_j}}\,dy ,
\end{align*}
by completing the square in the exponent and then substituting $y_j=x_j+\theta h_j$. So under $Q$ the
increments of $\tilde W$ are independent $\N(0,h_j)$, for every partition, and $\tilde W$ is a
$Q$-Brownian motion.
\end{proof}

The general theorem allows an adapted $\theta_t$, with
$Z_T=\exp(-\int_0^T\theta_t\,dW_t-\tfrac12\int_0^T\theta_t^2dt)$, provided $Z$ is a true martingale; Novikov's
condition $\E[\exp(\tfrac12\int_0^T\theta_t^2dt)]<\infty$ is enough \parencite{girsanov1960,karatzas1998}.
Its main use in this book is risk-neutral pricing. Under $P$ a stock follows
$dS=\mu S\,dt+\sigma S\,dW$. In a frictionless market with constant $r$, $\mu$ and $\sigma$ and the money
market account as numeraire, take $\theta=(\mu-r)/\sigma$, the market price of risk; under $Q$ the same
stock follows $dS=rS\,dt+\sigma S\,d\tilde W$. The drift changes and the volatility does not. Girsanov's
theorem supplies the change of drift; that this $Q$ is the measure to price with is an argument about
replication, which Chapter~5 makes.

Read in reverse, the theorem says that adding a drift is equivalent to reweighting:
$\E_P[f(W_T+\mu T)]=\E_P[f(W_T)e^{\mu W_T-\mu^2T/2}]$. With $f(x)=(x-0.5)^+$, $\mu=0.4$ and $T=1$, the exact
value is \girExact. Over \girPaths{} draws the direct estimate is \girDirect{} (standard error
\girDirectSE) and the reweighted estimate is \girWeighted{} (standard error \girWeightedSE).
The reweighted estimate is \girWeightedZ{} standard errors high. A miss this large in either direction
happens by chance in about one run in \girWeightedOdds, and one this far above in about one run in
\girWeightedOddsOne. Reweighting here costs variance, because the weights are spread out; the next subsection chooses the
weights so that they cut variance instead.

\subsection{Importance sampling}

Reweighting can also make simulation more efficient. Suppose we want the probability
$p=P(W_1>c)$ of an event in the far tail: a large loss, a deep default, or a three standard deviation
move in the five year yield, which on the pooled volatility of \canSigma{} points a year from \cref{sec3:qv} is a move
of about 2.5 points in a year. Plain simulation wastes almost every draw. Under the measure $Q$ of
\cref{thm3:girsanov} with $\theta=-c$, $W_1$ has mean $c$, half of the draws land in the event, and each
is reweighted by $dP/dQ=e^{-cW_1+c^2/2}$.

\begin{proposition}\label{prop3:is}
Let $Y\sim\N(c,1)$ and $\hat p=\mathbf 1\{Y>c\}e^{-cY+c^2/2}$. Then $\E[\hat p]=p=1-\Phi(c)$ and
\begin{equation}\label{eq3:isvar}
  \Var[\hat p] = e^{c^2}\bigl(1-\Phi(2c)\bigr)-p^2 .
\end{equation}
\end{proposition}

\begin{proof}
With $\phi$ the standard normal density, $\phi(y-c)e^{-cy+c^2/2}=\phi(y)$, so
$\E[\hat p]=\int_c^\infty\phi(y-c)e^{-cy+c^2/2}\,dy=\int_c^\infty\phi(y)\,dy=p$. For the second moment,
\begin{equation*}
  \E[\hat p^2]=\int_c^\infty\phi(y)e^{-cy+c^2/2}\,dy=e^{c^2/2}\int_c^\infty\phi(y)e^{-cy}\,dy .
\end{equation*}
Completing the square,
$\phi(y)e^{-cy}=e^{c^2/2}\phi(y+c)$, so the integral is $e^{c^2/2}(1-\Phi(2c))$ and
$\E[\hat p^2]=e^{c^2}(1-\Phi(2c))$.
\end{proof}

For $c=3$, $p=\isP\times10^{-4}$. The variance of one plain draw is $p(1-p)$, and by \cref{eq3:isvar} the
reweighted draw has about $1/\isRatio$ of that variance, so one reweighted draw is worth about \isRatio{}
plain ones. \Cref{fig3:girsanov}(b) shows the effect: none of 1{,}000 plain draws exceeds 3, so the
estimate is \isPlainSmall{} (in units of $10^{-4}$) with a sample standard error of \isPlainSmallSE. The
true standard error at this sample size is \isPlainSmallTrueSE, so the sample error bar presents a
useless answer as an exact one. By contrast 1{,}000
reweighted draws give \isSmall{} with a standard error of \isSmallSE. With 300{,}000 draws the two
estimates are \isPlainBig{} and \isBig, with standard errors \isPlainBigSE{} and \isBigSE.

\begin{figure}[t]
  \centering
  \includegraphics[width=\textwidth]{ch03/fig11_girsanov}
  \caption{(a) The law of $W_1$ under $P$ and under the tilted measure $Q$; the shaded area is
  $P(W_1>3)$. (b) Estimates of $P(W_1>3)$ by plain simulation and by importance sampling, with two
  standard errors; the dashed line is the exact value.}
  \label{fig3:girsanov}
\end{figure}

\FloatBarrier
% =====================================================================================
\section{Prices as expectations and as PDE solutions: Feynman-Kac}
\label{sec3:fk}

Many quantities in finance are discounted expectations of a function of a diffusion at a future date.
\textcite{kac1949}, building on Feynman's path integrals, showed that such expectations solve a
parabolic partial differential equation. That gives two independent ways to compute a price, averaging simulated payoffs or solving the
equation on a grid, and \cref{sec3:cn} does both for a call.

\begin{theorem}[Feynman-Kac]\label{thm3:fk}
Let $dX_s=a(X_s)\,ds+b(X_s)\,dW_s$ have a unique strong solution, let $r\ge0$ be a constant, and let
$u(t,x)$ be continuously differentiable once in $t$ and twice in $x$ on $[0,T]\times\mathbb R$ with
\begin{equation}\label{eq3:fkpde}
  u_t+a(x)u_x+\tfrac12b(x)^2u_{xx}-ru=0,\qquad u(T,x)=g(x).
\end{equation}
If $\E\int_t^T\bigl(b(X_s)u_x(s,X_s)\bigr)^2ds<\infty$ for the process started at $X_t=x$, then
\begin{equation}\label{eq3:fkrep}
  u(t,x)=\E\Bigl[e^{-r(T-t)}g(X_T)\Bigm| X_t=x\Bigr].
\end{equation}
\end{theorem}

\begin{proof}
Start the process at $X_t=x$ and let $Y_s=e^{-r(s-t)}u(s,X_s)$ for $s\in[t,T]$. By It\^o's lemma and the
product rule,
\begin{equation*}
  dY_s=e^{-r(s-t)}\Bigl(-ru+u_t+a\,u_x+\tfrac12b^2u_{xx}\Bigr)ds+e^{-r(s-t)}b\,u_x\,dW_s
       =e^{-r(s-t)}b(X_s)u_x(s,X_s)\,dW_s ,
\end{equation*}
because the bracket vanishes by \cref{eq3:fkpde}. So $Y_T-Y_t$ is an It\^o integral whose integrand is in
$L^2_{\mathrm{ad}}$ by assumption, and by \cref{thm3:itoext} it has mean zero. Hence
$u(t,x)=Y_t=\E[Y_T]=\E[e^{-r(T-t)}u(T,X_T)]=\E[e^{-r(T-t)}g(X_T)]$.
\end{proof}

\begin{remark}
The hypotheses can be weakened in two ways that matter below. The same proof works on an open state
interval, such as $(0,\infty)$ for a price, as long as the process stays in it. And $u$ need only be
smooth on $[0,T)\times\mathbb R$ and continuous up to $T$: apply the argument up to $T-\varepsilon$, where
the integrability condition must hold, and let $\varepsilon\downarrow0$. For the call below
$0\le u(t,s)\le s$, and since $e^{-rt}S_t$ is a $Q$-martingale closed by $e^{-rT}S_T$, the variables
$u(T-\varepsilon,S_{T-\varepsilon})$ are uniformly integrable. With the continuity of $u$ at $T$ they
converge in $L^1$ to the payoff, which justifies passing to the limit under the expectation.
\end{remark}

Two cases can be checked by hand. For Brownian motion ($a=0$, $b=1$, $r=0$) and $g(x)=x^2$, the PDE
$u_t+\tfrac12u_{xx}=0$ is solved by $u=x^2+(T-t)$, and indeed $\E[(x+W_{T-t})^2]=x^2+T-t$. For the
Ornstein-Uhlenbeck process with $g(x)=x^2$, \cref{eq3:oumoments} gives
$u(t,x)=\bigl(\theta+(x-\theta)e^{-\kappa\tau}\bigr)^2+\sigma^2(1-e^{-2\kappa\tau})/(2\kappa)$ with $\tau=T-t$.
Writing $m=\theta+(x-\theta)e^{-\kappa\tau}$, the derivatives are $u_\tau=-2\kappa m(x-\theta)e^{-\kappa\tau}+\sigma^2e^{-2\kappa\tau}$,
$u_x=2me^{-\kappa\tau}$ and $u_{xx}=2e^{-2\kappa\tau}$, so
$-u_\tau+\kappa(\theta-x)u_x+\tfrac12\sigma^2u_{xx}=2\kappa m(x-\theta)e^{-\kappa\tau}-2\kappa m(x-\theta)e^{-\kappa\tau}
-\sigma^2e^{-2\kappa\tau}+\sigma^2e^{-2\kappa\tau}=0$, using $u_t=-u_\tau$.

The case that matters most for this book is geometric Brownian motion under the measure $Q$ of
\cref{sec3:girsanov}, $a(s)=rs$ and $b(s)=\sigma s$, with the call payoff $g(s)=(s-K)^+$. Then
\cref{eq3:fkpde} is the Black-Scholes equation,
\begin{equation}\label{eq3:bs}
  V_t+rsV_s+\tfrac12\sigma^2s^2V_{ss}-rV=0,\qquad V(T,s)=(s-K)^+,
\end{equation}
and \cref{eq3:fkrep} says the price is the discounted $Q$-expectation of the payoff, once Chapter~5
has shown by replication that $Q$ is the measure to price with (the kink at
$s=K$ and the state space $(0,\infty)$ are covered by the remark after \cref{thm3:fk}). Chapter~5 derives
\cref{eq3:bs} a second way, by hedging, and solves it in closed form. Here we solve it numerically.

\subsection{A Crank-Nicolson solver}
\label{sec3:cn}

In the variable $x=\ln s$ the coefficients of \cref{eq3:bs} are constant: $a=r-\tfrac12\sigma^2$ and
$b^2=\sigma^2$. Put a uniform grid $x_0<\dots<x_N$ with spacing $\Delta x$ and run time backwards from
$T$ in steps of $\Delta\tau$. Replacing the $x$ derivatives by central differences gives, at each
interior node, the operator
\begin{equation*}
  (\mathcal Lu)_i=\Bigl(\frac{b^2}{2\Delta x^2}-\frac a{2\Delta x}\Bigr)u_{i-1}
  -\Bigl(\frac{b^2}{\Delta x^2}+r\Bigr)u_i+\Bigl(\frac{b^2}{2\Delta x^2}+\frac a{2\Delta x}\Bigr)u_{i+1},
\end{equation*}
and the Crank-Nicolson step averages the explicit and implicit schemes:
\begin{equation}\label{eq3:cn}
  \bigl(I-\tfrac12\Delta\tau\,\mathcal L\bigr)u^{k+1}=\bigl(I+\tfrac12\Delta\tau\,\mathcal L\bigr)u^k ,
\end{equation}
a tridiagonal system solved in $O(N)$ operations. For smooth data the scheme has error $O(\Delta x^2+\Delta\tau^2)$.
The kink in the call payoff spoils that at first, so the solver takes four fully implicit
half steps before switching to \cref{eq3:cn}, the smoothing device of \textcite{rannacher1984}, which
damps the oscillations the kink would otherwise cause. The grid in $x=\ln s$ is truncated at $s=1$ and
$s=1{,}000$, with boundary values $0$ at the bottom and $s-Ke^{-r\tau}$ at the top.

With $K=100$, $r=3\%$, $\sigma=20\%$ and one year to expiry, the solver gives \fkPDE{} at $s=100$,
equal to the closed form \fkBS{} to four decimals. A simulation of \fkPaths{} terminal values under $Q$,
drawn as 500{,}000 antithetic pairs, gives \fkMC{} with a standard error of \fkMCse. For the Ornstein-Uhlenbeck second moment with
$\kappa=1.5$, $\theta=0.2$, $\sigma=0.5$, $T=2$ and $x=1$, the solver gives \fkOUPDE, the closed form
\fkOUExact{} and simulation \fkOUMC. Panel (b) of \cref{fig3:fk} shows second order convergence for this grid, boundary and
time-stepping set-up, with fitted slopes of \fkSlopeCall{} for the call and \fkSlopeOU{} for the moment;
with rougher data or other boundaries the order would have to be checked again.

\begin{figure}[t]
  \centering
  \includegraphics[width=\textwidth]{ch03/fig13_feynman_kac}
  \caption{(a) The value of a one year call computed by solving \cref{eq3:bs} with Crank-Nicolson, by
  the closed form and by simulation. (b) Error of the PDE solution against the number of space steps,
  for the call and for $\E[X_T^2]$ under the Ornstein-Uhlenbeck process.}
  \label{fig3:fk}
\end{figure}

\FloatBarrier
% =====================================================================================
\section{Case study: is the five year yield mean reverting?}
\label{sec3:case}

Mean reversion is often assumed for rates, spreads, cap rates and valuation multiples, and a fitted
half-life is quoted as if it were a measured property of the market. This case study fits
\cref{eq3:ou} to the month-end five year Government of Canada benchmark yield from \caseStart{} to
\caseEnd, \caseN{} observations from the archived snapshot used in Chapter~1, and asks how much the fit can
be trusted. The snapshot starts in April 2009 because it pairs the yield with the policy rate
target, whose series in the file begins then; the sample ends at the last complete month-end in the
data. Refitting on the longer Valet series of \cref{sec3:qv}, \caseLongN{} month-ends from
\caseLongStart, gives $\hat\kappa=\caseLongKappa$, a half-life of \caseLongHalf{} years and a
Dickey-Fuller $p$-value of \caseLongADFp, lower than on the shorter sample but still far from rejecting a random walk, and the checks below
come out the same way.

Regressing each month-end yield on the previous one gives $\hat\phi=\casePhi$. By \cref{cor3:ar1} this
means $\hat\kappa=\caseKappa$ a year, a long-run level of $\hat\theta=\caseTheta\%$, a volatility of
$\hat\sigma=\caseSigma$ points a year and a half-life of \caseHalf{} years. The last observation is
\caseLast\%. Panel (a) of \cref{fig3:yieldou} shows the series against the fitted level and the
stationary band $\hat\theta\pm1.96\,\hat\sigma/\sqrt{2\hat\kappa}$, which reaches \caseBand{} points either side of $\hat\theta$ because the stationary standard deviation is \caseSdInf.

Three checks make the numbers look much less solid. First, the Dickey-Fuller statistic (with a
constant and no lagged differences) for the null hypothesis $\phi=1$, a random walk with no reversion
at all, is \caseADF, against a 5\% critical value of \caseADFcrit{} \parencite{dickey1979}; the
$p$-value is \caseADFp, so the data cannot reject a random walk. With \caseN{} monthly points the test has little power
against a half-life of three years, so failing to reject is weak evidence either way. Second, a parametric bootstrap that
simulates \caseBoot{} samples of the same length from the fitted model and refits each one finds an
average $\hat\kappa$ of \caseKappaBootMean, \caseBiasRatio{} times the value used to simulate (panel b).
In simulations calibrated to this fitted model the estimator is badly biased upwards at this sample
size; the ratio belongs to this experiment and would differ at another sample size or speed. Third, under the
fitted model, 90\% of the half-lives estimated from samples this long lie between \caseHalfLo{} and
\caseHalfHi{} years. That is the spread of the estimator in simulation, not a confidence interval for
the true half-life, and it shows that even if the model were right the estimate would be imprecise.

Seventeen years of monthly data are consistent with slow mean reversion but do not establish it. A
model that relies on a fitted half-life of about three years for, say, the expected holding period of
a relative value trade is relying on a number the data barely support. Longer samples help, and so
does pooling across related series. An approximate correction for the bias uses the first-order formula
of \textcite{kendall1954}, but with a coefficient this close to one the corrected value depends heavily
on the approximation. Here Kendall's correction with $n=\caseNK$, the number of
regression pairs, moves $\hat\phi$ from \casePhi{} to \casePhiKendall{} (with $n=\caseN$, the number of
month-ends, it gives \casePhiKendallAlt), which puts $\hat\kappa$ at essentially zero either way: the fitted speed of
mean reversion is not robust.

\begin{figure}[t]
  \centering
  \includegraphics[width=\textwidth]{ch03/fig10_yield_ou}
  \caption{(a) Month-end five year Government of Canada yield with the fitted long-run level and
  stationary band. (b) Estimates of $\kappa$ from \caseBoot{} samples simulated from the fitted model; the
  estimator is biased upwards.}
  \label{fig3:yieldou}
\end{figure}

\FloatBarrier
% =====================================================================================
\section{Summary}
\label{sec3:summary}

A random walk with small fast steps becomes Brownian motion, a continuous process with independent
Gaussian increments that can be built by repeatedly filling in midpoints. Its paths are rough at every
scale. Their total variation is infinite, but their quadratic variation converges to $T$ along the
refining partitions of \cref{thm3:qv}, and this, written $(dW)^2=dt$, is where nearly every correction term
that follows comes from. The realised quadratic variation of the five year Canadian yield gives an empirical
comparison, though it cannot show that the yield is a Brownian motion with constant volatility: its median realised volatility is \canRvDailyMed{} points from daily changes and \canRvMonthlyMed{} from
month-end changes, a gap no larger than twelve observations a year would produce by chance, although
its volatility changes from year to year.

The It\^o integral is defined with the integrand fixed at the start of each interval, which is how a
trading position works. Its isometry controls the second moment, while adaptedness and the
independence of Brownian increments give the martingale property. It\^o's lemma adds the term $\tfrac12\sigma^2f_{xx}\,dt$ to the ordinary chain rule; for a five year
zero coupon bond it gives the convexity term in the expected return. With it,
geometric Brownian motion and the Ornstein-Uhlenbeck process can be solved in closed form, and the
first shows why the median of a volatile asset can fall while its mean rises.

Martingales and optional stopping answer questions about barriers, such as which of two levels is hit
first. Girsanov's theorem shows that reweighting paths by an exponential martingale adds a drift to a
Brownian motion and leaves its volatility alone, which is the basis of risk-neutral pricing and of importance sampling. The Feynman-Kac
formula identifies discounted expectations with solutions of parabolic PDEs, so a price can be
computed by simulation or on a grid. The case study shows that fitting these models is a
separate problem from solving them: in simulations calibrated to the five year yield, the fitted speed
of mean reversion averages more than twice the value used to generate the data.

\FloatBarrier
% =====================================================================================
\section{Exercises}
\label{sec3:exercises}

\begin{enumerate}[label=\textbf{3.\arabic*}, leftmargin=*, ref=3.\arabic*]
  \item\label{ex3:one} Using $\Cov(W_s,W_t)=\min(s,t)$, find $\Var(W_1+W_2)$ and $\Cov(W_2,W_3-W_1)$.
  \item\label{ex3:two} Apply It\^o's lemma to $W_t^4$ and take expectations to show $\E[W_t^4]=3t^2$. What is
  $\E[W_2^4]$?
  \item\label{ex3:three} A stock follows geometric Brownian motion with $S_0=100$, $\mu=8\%$ and
  $\sigma=25\%$. Find the mean and median of $S_2$ and $P(S_2<100)$.
  \item\label{ex3:four} An Ornstein-Uhlenbeck spread has $\kappa=0.5$ a year and $\sigma=2\%$ a year. Find
  its half-life, its stationary standard deviation and the correlation between values one year apart in
  the stationary state.
  \item\label{ex3:five} Find $\Var\bigl(\int_0^3t\,dW_t\bigr)$ and $\Var\bigl(\int_0^3W_t\,dt\bigr)$.
  \item\label{ex3:six} For $X_t=\mu t+\sigma W_t$ with $\mu=5\%$ and $\sigma=20\%$, find the probability that
  $X$ reaches $0.1$ before $-0.2$, and the expected time until it reaches one of them. Repeat with $\mu=0$.
  \item\label{ex3:seven} Let $Z_T=\exp(\theta W_T-\tfrac12\theta^2T)$. Use \cref{thm3:girsanov} to find
  $\E[Z_TW_T]$ for $\theta=0.5$ and $T=2$ without computing an integral.
  \item\label{ex3:eight} Solve $u_t+\tfrac12u_{xx}=0$ with $u(T,x)=x^2$ and evaluate $u(0,1)$ when $T=2$.
  Then use \cref{thm3:fk} to find $\E[X_T^2]$ for $dX=0.1\,dt+0.3\,dW$, $X_0=1$, $T=2$.
  \item\label{ex3:nine} Show that $\sum_i(W_{m_i}-W_{t_i})(W_{t_{i+1}}-W_{t_i})$, with $m_i$ the midpoints,
  has mean $T/2$, and explain why this makes the midpoint rule unsuitable for modelling trading gains.
\end{enumerate}

\subsection*{Answers}

\begin{enumerate}[label=\textbf{3.\arabic*}, leftmargin=*]
  \item $\Var(W_1+W_2)=1+2+2\min(1,2)=\exOneVar$. The covariance is
        $\Cov(W_2,W_3-W_1)=\min(2,3)-\min(2,1)=\exOneCov$.
  \item $d(W^4)=4W^3dW+6W^2dt$, so $\E[W_t^4]=6\int_0^t\E[W_s^2]\,ds=6\int_0^ts\,ds=3t^2$ and
  $\E[W_2^4]=\exTwo$.
  \item $\E[S_2]=100e^{0.16}=\exThreeMean$. The median is $100e^{(0.08-0.03125)\cdot2}=\exThreeMedian$.
  $P(S_2<100)=\Phi\bigl(-0.04875\sqrt2/0.25\bigr)=\exThreeLoss\%$.
  \item Half-life $\ln2/0.5=\exFourHalf$ years; stationary standard deviation
  $0.02/\sqrt{2\cdot0.5}=\exFourSd\%$; correlation $e^{-0.5}=\exFourAcf$, from \cref{eq3:oumoments} as
  $s\to\infty$ with $t-s=1$.
  \item By the isometry, $\int_0^3t^2dt=\exFive$. By \cref{ex3:parts}, $\int_0^3W_t\,dt=\int_0^3(3-t)\,dW_t$,
  with variance $\int_0^3(3-t)^2dt=\exFive$ as well.
  \item With $k=2\mu/\sigma^2=2.5$, \cref{eq3:exit} gives $p=\exSixP\%$ and $\E[\tau]=\exSixTau$ years. With $\mu=0$, $p=0.2/0.3=\exSixPZero\%$ and $\E[\tau]=0.1\cdot0.2/0.04=\exSixTauZero$ years.
  \item $Z_T$ is the density of a measure under which $W_t-\theta t$ is a Brownian motion (\cref{thm3:girsanov}
  with $-\theta$). So $\E[Z_TW_T]$ is the mean of $W_T$ under that measure, $\theta T=\exSeven$.
  \item $u(t,x)=x^2+(T-t)$, so $u(0,1)=\exEight$. For the drifted process the PDE of \cref{thm3:fk}
  is $u_t+0.1u_x+0.045u_{xx}=0$ with $u(T,x)=x^2$, solved by $u(t,x)=(x+0.1(T-t))^2+0.09(T-t)$, so
  $\E[X_T^2]=u(0,1)=(1+0.1\cdot2)^2+0.3^2\cdot2=\exEightB$.
  \item Write the term as $\sum_iA_i(A_i+B_i)$ as in the proof of \cref{prop3:strat}; its mean is
  $\sum_i\E[A_i^2]=\sum_ih_i/2=T/2$. A strategy that could choose its position over $[t_i,t_{i+1})$
  using $W_{m_i}$ would be using information from the middle of the interval, after the position has to
  be set, so the extra $T/2$ is a gain from seeing the future.
\end{enumerate}

\FloatBarrier
% =====================================================================================
\section{Notes, sources and credits}
\label{sec3:credits}

\paragraph{Sources}
The construction of Brownian motion follows the midpoint method in \textcite{morters2010}; the
continuity correction for discrete monitoring is from \textcite{broadie1997}. It\^o's integral goes back to
\textcite{ito1944}, Girsanov's theorem to \textcite{girsanov1960} and the Feynman-Kac formula to
\textcite{kac1949}. \textcite{karatzas1998} and \textcite{oksendal2003} give complete proofs of the
results stated here with outlines; \textcite{steele2001}, \textcite{shreve2004} and \textcite{wilmott2006}
are gentler treatments aimed at finance. On simulation, \textcite{kloeden1992} is the standard reference
for convergence orders, \textcite{higham2001} a short practical introduction and
\textcite{glasserman2003} the reference for importance sampling. The estimation bias in the case study is
discussed by \textcite{kendall1954}, and the unit root test is that of \textcite{dickey1979}. The
convergence acceleration used in the R cross-check to evaluate $\zeta(\tfrac12)$ is from
\textcite{cohen2000}.

\paragraph{Open source projects}
\begin{sloppypar}
The computed results come from eight scripts in \texttt{code/ch03/} (\texttt{random\_walks},
\texttt{brownian\_qv}, \texttt{ito\_integral}, \texttt{sde\_models}, \texttt{martingales\_girsanov},
\texttt{feynman\_kac}, \texttt{canadian\_rates} and \texttt{exercises}) and the short listing in
\cref{sec3:sde}. Each script draws from its own NumPy stream, seeded with the pair (303, $k$) for
script $k$, and the listing uses seed 1. The tests in \texttt{code/ch03/test\_*.py} check the scripts
against closed forms and simulation, and the R and Octave implementations recompute the closed forms
by other methods. The code uses NumPy \parencite{numpy} for arrays, arithmetic and random number generation,
SciPy \parencite{scipy} for special functions and the normal law, Matplotlib \parencite{matplotlib}
for the figures and statsmodels \parencite{statsmodels} to compute the Dickey-Fuller test reported in
\cref{sec3:case}. The \texttt{stochastic} \parencite{stochasticpy} and \texttt{sdeint} \parencite{sdeint} packages offer independent implementations of
these processes and integrators that a reader can use to check the simulations; the numbers in this
chapter come only from its own scripts. The notebooks of \textcite{cantaro86} and \textcite{py4fi} contain
treatments of Brownian motion, the Ornstein-Uhlenbeck process, Feynman-Kac and finite difference
methods, and the QuantLib library \parencite{quantlib} provides production implementations of these
processes and of finite difference solvers; nothing in this chapter calls it. The R and Octave
cross-checks run on base R and on GNU Octave \parencite{octave}. The versions cited are those pinned in
\texttt{code/requirements-lock.txt} and used to compute every number here, NumPy 2.2.6, SciPy 1.18.1,
Matplotlib 3.11.2 and statsmodels 0.15.0, and the Octave cross-check ran on GNU Octave 9.4.0. The full list, with links, is in
\texttt{CREDITS.md}.
\end{sloppypar}

\paragraph{Data}
The case study uses Bank of Canada series \texttt{BD.CDN.5YR.DQ.YLD} through the MIT-licensed snapshot
in \textcite{khan2026policy}, cut at the September 2026 month-end; the snapshot starts in April 2009, where its policy rate
series begins. The realised-variance, reflection, convexity and importance-sampling examples use the
same yield series from January 2001, pulled from the Bank of Canada Valet interface (accessed 19
October 2026) and archived as \texttt{data/goc\_bills\_bonds.csv}.

\printbibliography[heading=subbibliography,title={References for Chapter \thechapter}]
\end{refsection}
